\documentclass{jsarticle}
\usepackage[dvipdfmx,final]{graphicx}
\usepackage{amsmath}
\usepackage{amssymb}
\begin{document}
\title{}
\author{}
\date{}
\maketitle
   \newcommand{\variint}{%
	\begin{picture}(15,30)
		\put(0,0){
			$ \iint $}
		\put(0,5){\line(15,0){15}}
	\end{picture} %
}

\newcommand{\variiint}{%
	\begin{picture}(15,15)
		\put(0,0){
			$ \iiint $}
		\put(0,5){\line(15,0){20}}
	\end{picture} %
}



 \newcommand{\alearletter}{%
	 \begin{picture}(10,20)
		 \put(0,0){
			 $ \square $}
		 \put(0,0){\line(1,1){10}}
	 \end{picture} %
	 }

 \newcommand{\secproduct}{%
	 \begin{picture}(10,20)
		 \put(0,0){
			 $ \nabla $}
		 \put(5,0){+}
	 \end{picture} %
	 }
 \newcommand{\talot}{%
	 \begin{picture}(15,20)
		 \put(0,0){
			 $ \text{\scalebox{2.0}[2]{$\nabla$}}$}
		 \put(5,3){
			 $ \Delta $}
	 \end{picture} %
	 }
 \newcommand{\hazerdnabla}{%
	 \begin{picture}(15,20)
		 \put(0,0){
			 $ \text{\scalebox{2.0}[2]{$\nabla$}}$}
		 \put(5,3){
			 $\text{Y}  $}
	 \end{picture} %
	 }

 \newcommand{\securebox}{%
	 \begin{picture}(15,20)
		 \put(0,0){
			 $ \text{\scalebox{2.0}[2]{$\square$}}$}
		 \put(5,3){
			 $\text{Y}  $}
	 \end{picture} %
	 }


\newcommand{\timebow}{%
	\begin{picture}(5,5)
		\put(0,0){$\text{\rotatebox{90}{$\bowtie$}}$}
	\end{picture} %
	}


 \newcommand{\flowerletter}{%
	 \begin{picture}(10,20)
		 \put(0,0){$\text{\scalebox{1.0}[1]{$\square$}}$}
		 \put(0,0){$\text{\scalebox{1.0}[2]{$\int$}}$}
	 \end{picture} %
	 }
\newcommand{\squareiint}{%
	\begin{picture}(15,30)
		\put(0,0){
			$ \iint $}
		\put(0,5){$\text{\scalebox{2.0}[1]{$\square$}}$}
	\end{picture} %
}

\newcommand{\squareint}{%
	\begin{picture}(15,30)
		\put(0,0){
			$ \int $}
		\put(0,5){$\text{\scalebox{2.0}[1]{$\square$}}$}
	\end{picture} %
}


   \newcommand{\dazanier}{%
	\begin{picture}(12,12)
		\put(0,0){
			$ \bigsqcup $}
		\put(4,3){$\text{\scalebox{1.0}[1]{$\top$}}$}
		\put(4,3){$\text{\scalebox{1.0}[1]{$\bot$}}$}
	\end{picture} %
}
\newcommand{\innabla}{%
	\begin{picture}(13,13)
		\put(0,0){
			$ \nabla $}
		\put(0,0){$\text{\scalebox{2.0}[2]{$\nabla$}}$}
	\end{picture} %
}

\newcommand{\inndelta}{%
	\begin{picture}(20,20)
		\put(0,0){
			$ \Delta $}
		\put(0,0){$\text{\scalebox{2.0}[2]{$\Delta$}}$}
	\end{picture} %
}

\newcommand{\starnabla}{%
	\begin{picture}(20,20)
		\put(3,5){
			$ \Delta $}
		\put(0,0){$\text{\scalebox{2.0}[2]{$\nabla$}}$}
	\end{picture} %
}


  \newcommand{\whitewithbodercircle}
  {\text{\textcircled{\phantom{\textbullet}}}
  \kern-0.9em\text{\textcircled{\phantom{}}}}
     

 
  \newcommand{\whitewithblack}{\text{\textcircled{\textbullet}}}
  \newcommand{\whitewithblacksquare}{\text{\fbox{\text{\rule{0.5em}
  {0.5em}}}}}
  \newcommand{\whitewithwhitesquare}{\text{\fbox{\text{\phantom
  {\rule{0.5em}{0.5em}}}}}}
 \newcommand{\whitewithbodercircleiiint}{%
	 \begin{picture}(20,20)
		 \put(0,0){\scalebox{2.0}[2]{$\whitewithbodercircle$}}
		 \put(0,0){\scalebox{2.0}[2]{$\iiint$}}
	 \end{picture}%
	 }
\newcommand{\whitewithbodercircleint}{%
	 \begin{picture}(20,20)
		 \put(2,2){\scalebox{1.0}[1]{$\whitewithbodercircle$}}
		 \put(2,2){\scalebox{1.0}[2]{$\int$}}
	 \end{picture}%
	 }
\newcommand{\whitewithbodercircleiint}{%
	 \begin{picture}(20,20)
		 \put(0,0){\scalebox{1.0}[1]{$\whitewithbodercircle$}}
		 \put(2,2){\scalebox{1.0}[2]{$\iint$}}
	 \end{picture}%
	 }
 
 \newcommand{\whitewithbodercirclevarint}{%
	 \begin{picture}(13,13)
		 \put(0,0){\scalebox{2.0}[2]{$\whitewithbodercircle$}}
		 \put(0,0){\scalebox{2.0}[2]{$\varint$}}
	 \end{picture}%
	 }
 
 \newcommand{\whitewithbodercirclevariiint}{%
	 \begin{picture}(13,13)
		 \put(0,0){\scalebox{2.0}[2]{$\whitewithbodercircle$}}
		 \put(0,0){\scalebox{2.0}[2]{$\variiint$}}
	 \end{picture}%
	 }
 
 \newcommand{\average}{%
	  \begin{picture}(30,45)
		  \put(0,0){\scalebox{2.0}[2]{$\square$}}
		  \put(4,4){$\setminus$}
	  \end{picture}%
  } 
  \newcommand{\smallaverage}{%
	  \begin{picture}(30,45)
		  \put(0,0){\scalebox{1.0}[1]{$\square$}}
		  \put(0,0){$\setminus$}
	  \end{picture}%
  } 
 
  \newcommand{\whitewithboder}{%
	  \begin{picture}(15,30)
		  \put(0,0){\scalebox{2.0}[2]{$\square$}}
		  \put(4,4){$\square$}
	  \end{picture}%
  } 
  \newcommand{\whitewithnabla}{%
	  \begin{picture}(15,30)
		  \put(0,0){\scalebox{2.0}[2]{$\nabla$}}
		  \put(3,3){$\nabla$}
	  \end{picture}%
  } 



  Open integral have belong with open integral fundement gravity in
  delanversian element own deprivation.
  $$\talot = \oint_{D} M(\square) d\square $$
  And, this gravity equal with fundemental group.
  $$ \oint_M \pi(\chi,x) = \oint_M [i\pi(\chi,x),f(x)] $$
  $$ M(\square) = x^y + y^x + z^a + u^b + v^c = 0 $$
  Kalavi-Yau manifold.
  $$ {1 \over y} + {1 \over x} + {1 \over a} + {1 \over b} + {1 \over c} = 1 $$
  $$ \talot = \oint_D M(\square)dm, T = \Gamma^{'}(\gamma)dx_m $$
  $$ \square = 2(\sin(ix \log x) + i\cos(ix \log x)) $$
  Circumstance have with gravity equation.
  $$ = {d \over {d \gamma}} \Gamma $$
  $$ M = [i\pi(\chi,x),f(x)] $$
  $$ \square = -{{16\pi G}\over {c^4}}T^{\mu\nu} $$
  Gamma function in partial gravity of deprivation.
  $$ = \kappa T^{\mu\nu} $$
  These equation is concluded with general relativity theory.

  D-brane are also constructed from Thurston Perelman manifold.
  More also, this equation is constructed with quantum formula.
  $$ \int E^{'}(\sigma)d\sigma = \nabla_{i}\nabla_{j} \bigoplus
  (H(\sigma) \otimes K(\sigma))\nabla \eta d \eta $$
  $$ \sigma = \int{(h\nu)^{\nabla}}^{\oplus L}d\Psi $$

  Secure product is own have with quantum level of gravity equation.
  $$ \secproduct {(\square(\nabla \psi)^{\nabla}}^{\oplus L} 
  = \int \square^{'}(\nabla \psi)dx_m = \boxplus \Psi $$
  $$ x \boxtimes y = \bigoplus \nabla w $$
  $$ =(\boxtimes x)^{x + y} $$
  Projection of equation have with box element and category theory.
  $$ (\bigoplus \nabla w) \boxtimes (\bigoplus \nabla w) $$
  $$ = \bigoplus{(i\hbar^{\nabla})}^{\oplus L} $$
  Quantum level of space ideality equation also have with factor equation.
  $$ (\bigoplus {(i\hbar^{\nabla})}^{\oplus L} + m) 
  (\bigoplus {(i\hbar^{\nabla})}^{\oplus L} + n)
  = {{L^{m+1}} \over {m + 1}} = \int (x-1)^{t-1}\cdot t^{x-1}dt $$
  $$ = \beta(p,q) $$
  And these equation conclude with beta function.
  $$ \square^{{{x+y} \over 2} - \sqrt{x\cdot y}} = \square^{\ll o} $$
  $$ \square \Psi \boxtimes \square \Psi $$
  Dalanversian equation of zeta function own have with average equation.
  $$ \boxtimes = \dazanier, \dazanier^{{1 \over 2}} = \secproduct $$
  Tunnel daiord of equation is belog with Jones manifold, and
  this equation also have belong with zeta function of value in
  deprivation of element.
  $$ \flowerletter {{Z(\zeta)} \over {h\nu}}dx_{\zeta} $$
  Under equations comment with Euler function of circumstance formula
  $$ = \flowerletter {{Z(\zeta)} \over {\log x}}dx_{\zeta} $$
  $$ S = \pi \oint ||r^2||dr $$
  $$ = \int e^{-x^2 - y^2} dxdy \cdot \oint ||r^2||dr $$
  $$ = 2 \pi S $$
  $$ e^{2\pi r} = 1 $$
  After all, zeta function also mention to build with quantum element
  in circle function resolved from Gauss function.

  $$ \Gamma |: r \to \chi = \hazerdnabla \to \securebox \to [
	  \Delta, \nabla, d,\partial,\delta,dx_m] $$
  $$ \to Y|:m \to n $$
  Gamma function also construct with deprivation of element.
  $$ {d \over {df}}F(x,y) = \bigoplus[{{\pi(\chi,x)} \over {x\log x}}dx_m
  + i{{\pi(\chi,x)} \over {x\log x}}dy_m][I_m]$$
  Higgs function of average equation equal with Cauchy function of Euler
  equation.
  $$ \secproduct [\int \pi(\chi,x)dx_m + \int {}^{N}\pi(\chi,x)dy_m][dI_m] $$
  $$ \int d \timebow = \int f(x)d(x\log x) $$
  $$ = F $$
  $$ \timebow = {f \over {x\log x}} $$
  After all, time of deprivate value is logment element.

  Jones manifold estimate with space ideality from partial gamma function
  of integral formula to Higgs field dependent with quata equation.
  $$ {{e^{f} + e^{-f}} \over {e^{f} - e^{-f}}} = {{\int \Gamma^{'}(\gamma)
  dx_m} \over {m(x)}} $$
  $$ {d \over {d f}}F(x,y) = m(x,y) $$
  This equation also constructed with fundamental group of time scale value.
  $$ \square_{k = prime}^{\infty} Z^{\ll D{\timebow}}(\zeta)[I_m] $$
  $$ = \pi(\timebow\cdot \bar{\timebow},m) $$
  Mebius formula is included with triple varint integral equation project 
  with seed of pole annouce.
  $$ {}^{t}\scalebox{1.0}[2]{\variiint}_{D\chi}[\pi(\chi,x) |: x \to
  2, |:y \to \infty][dI_m]$$
  $$ = \scalebox{1.0}[2]{\variiint}[{{{\timebow_{x}\cdot \bar{\timebow_m}}
  - {\bar{\timebow_y}\cdot \timebow_n}}\over {t - t_1}}]d\timebow_m $$
  These equation equal with beta function.
  $$ \beta(p,q) = \secproduct[\timebow_m|_{x,y} \cdot \bar{\timebow}|_{x,y}
  ]\times [\sigma(\timebow)] $$
  Thuston Perelman manifold are endeavor with being constructed from being
  stuggled to being mixin with mebius dimension.
  $$ E = K(\sigma) \otimes H(\sigma) = \beta^{-1}(x)x\beta(x) $$
  Under equations are projection of fundament group from cross of operate to
  dalanversian equation being inervibled with gamma function of partial
  deprivation being deconstructed a gamma value into cauchy of zeta integral
  equation.
  $$ \innabla |:\chi \to \nabla_{i}\nabla_{j}\Gamma^{'}(\gamma)d\gamma_{r_m} $$
  $$ \to \starnabla \square \Psi $$
  $$ \timebow(H\Psi_{D\chi})^{\ll D\psi} $$
  $$ = \oint {{Z^{'}(\zeta)} \over {2\pi i}}[dN] $$
  $$ ={{[f(x)g^{'}(x) - g(x)f^{'}(x)] - [F(x)G(x)]} \over {x \log x}} $$
  $$ = {{i\pi(\chi,x)} \over {\log x}} = \Delta f - \nabla g $$
  $$ = ({d \over {d f}Fg})\cdot (g(f)) $$
  $$ {\partial \over {\partial f}}F(x,y) = (F^{f}(x,y)f^{'}(x,y) $$
  $$ = {\partial \over {\partial f}}F(g(x,y)) $$ 
  Under equation is constructed with time scale value of pair from star of
  quantum level to Jone manifold being explained with flow of time from
  universe to other dimension.
  $$ \whitewithbodercircleiint^{N} (\timebow \cdot \bar{\timebow})_{D}
  = \inndelta_{D\chi}{(\ast^{\nabla})}^{\ll D}$$
  Particle operate emelite with fundament group of varint integral of
  quantum equation.
  $$ {}^{t}\scalebox{1.0}[2]{\variint}_{D} 
  \pi(\chi,x)_m = \ll |\emptyset | \alearletter| \gg $$

  $$ y = e^{n}, \text{if n select with prime number, then y is productivity 
  number,}$$ 
  neipia number times n is shanon entropy belong with.
  moreover, this times escourt into even number. 
  $$ u = {e^{n} \over {\log 2}}, \text{if $e^{n}$ select with $n = x \log x$ 
  , then u sensevility have with zeta function} $$
  After all, $ x^{{1 \over 2} + iy} = e^{x \log x} $ is zeta function.
  Zeta function also have with cauchy law equation.


  $$ e^{f} + e^{-f} \ge e^{f} - e^{-f} $$
  $$  (e^{f} + e^{-f} \ge e^{f} - e^{-f})^{'} $$
  $$ = e^{f} - e^{-f} \le e^{f} + e^{-f} $$
  $$ R^{'}_{ij} = - R_{ij} $$
  $$ = -2 R_{ij} \ge 0 $$

  極限値は、0であり、
  $$ R^{'}_{ij} \ge 0 $$
  である。
  $$ \bigoplus{(i \hbar^{\nabla})}^{\oplus L} = e^{-f} $$
  $$ {L^{\nabla}}^{\oplus L} - {{L^{\nabla}}^{\oplus L}}^{-1} 
  \le {L^{\nabla}}^{\oplus L} + {{L^{\nabla}}^{\oplus L}}^{-1} \le
  e^{f} - e^{-f} \le e^{-f} + e^{f} $$
  $$ \left({L^{\nabla}}^{\oplus L} + {{L^{\nabla}}^{\oplus L}}^{-1}\right)^{df}
  = \sin \theta \ge {\sqrt{3} \over 2} $$
  $$ \int C dx_m = \int(\int {1 \over x^s}dx - \log x)\mathrm{dvol} $$
  $$ = -2 R_{ij} = 2 {d \over df}F $$
  $$ {d \over df}F + \int C dx_m \ge {d \over df}F - \int C dx_m $$
  $$ = e^{x \log x} + {1 \over e^{x \log x}} \le e^{x \log x} - {1 \over e^{x \log x}} $$
  $$ \ge \sin \theta $$
  $$ \ge {\sqrt{3} \over 2} $$
  $$ \cong \cos \theta \ge {1 \over 2}$$
  $$ {d \over df}F + \int C dx_m \ge {d \over df}F - \int C dx_m $$
  $$ = 2(\cos(i x\log x) - i\sin(i x\log x)) $$
  $$ = 2(-\sin(x \log x) + i\cos(x \log x) $$
  $$ = ({1 \over  2}i - {{\sqrt{3} \over 2})} $$
  $$ = e^{-\theta} $$
  $$ = \cos \theta + i \sin \theta $$
  $$ = {{1 \over 2}} + i{{\sqrt{3} \over 2}}$$
  $$ = e^{i\theta} $$
  $$ \int \int{1 \over (x \log x)^{2}}dx_m \ge {1 \over 2}i, \int \int{1 \over (y \log y)^{1 \over 2}}dy_m
  \ge {1 \over 2} $$
  $$ = e^{-\theta} + e^{i\theta} $$
  $$ $$

  オイラーの定数を多様体積分した値は、2倍の値のゼータ関数の指数作用であり、
  $$ \square^{x + iy} = \square^{i} = 2\left({d \over df}F + \int C dx_m\right) = \left({d \over df}\int \int
  {1 \over {(x \log x)^2}}dx_m + {d \over df}\int\int{1 \over {(y \log y)^{1 \over 2}}}dy_m\right) $$
  正常な細胞のヒッグス場のエネルギーからオイラーの定数のガンマ関数の大域的積分多様体
   のエネルギーを差し引いたエネルギーであり、宇宙から異次元
  への角度の数値でもあり、このオイラーの定数のエネルギーバランスが、永遠の生命エネルギーの
  数値へのバランスであり、彩さんの過去のレビューからの導きでもあり、
  私の集大成であり、広島大学医学部からの導きであることは、
  私の論文からも、一目瞭然であり、富山大学への感謝である。
   $$\beta(p,q) = \int||\sin^{2} \theta||d\tau $$
   $$ (\sin \theta \ge {\sqrt{3} \over 2}) $$
   の$\theta \ge 60$度の範囲である。
   以下は、この値が生成される元の式たちである。
   $$ \int \int{1 \over (x \log x)^2}dx_m \ge {1 \over 2}i, \int \int{1 \over (y \log y)^{1 \over 2}}dy_m
   \ge {1 \over 2} $$
   $$ \int \int{1 \over (x \log x)^2}dx_m \ge {1 \over 2}i $$
   $$ \ge {\sin \theta } \ge {\sqrt{3} \over 2} $$
    $\sin 30^{\circ} = \cos 60^{\circ}$. これを一般座標変換すると、 
   $$ \int \int{1 \over (y \log y)^{1 \over 2}}dy_m
   \ge {1 \over 2} $$
   $$ \ge {\cos \theta} \ge {1 \over 2} $$
   $$ \int(\int {1 \over x^s}dx - \log x)\mathrm{dvol} \ge C \ge \cos 30^{\circ} \ge {1 \over 2} $$
   $$ \int {1 \over x^s}dx - \log = C $$
   $C = 0.571....$であり、
   $$ \int C dx_m = \int||\sin^{2} \theta||d\tau \ge {\sqrt{3} \over 2} $$
   $$ = \int e^{2\sin \theta \cos \theta}\int \sin \theta \cos \theta d\theta $$
   である。
   $$ {d \over df} \int C dx_m = {{\sqrt{3} \over 2}^{{\sqrt{3} \over 2}}}^{'} $$
   $$ = 1 $$
   これは、以下の式から、この式の構造がどうなっているかが分かる。

   $$ {L^{\nabla}}^{\oplus L} - {{L^{\nabla}}^{\oplus L}}^{-1} 
  \le {L^{\nabla}}^{\oplus L} + {{L^{\nabla}}^{\oplus L}}^{-1} \le
  e^{f} - e^{-f} \le e^{-f} + e^{f} $$
  $$ \left({L^{\nabla}}^{\oplus L} + {{L^{\nabla}}^{\oplus L}}^{-1}\right)^{df}
  = \sin \theta \ge {\sqrt{3} \over 2} $$

\section{Hilbert manifold in Mebius space \\
        this element of Zeta function on integrate of fields}

 Hilbert manifold equals with Von Neumann manifold, and this fields is concluded with Glassman manifold,
the manifold is dualty of twister into created of surface built.
These fields is used for relativity theorem and quantum physics.
$$ ||ds^2|| = \lim_{x \to \infty} [\delta(x) \int \int \int \pi \left( \sum_{k=0}^{\infty}
		{{}^n\sqrt{p},x \over n} \right)
^{1 \over 2}d \tau]^{\mu\nu} $$
  This norm is component of fields in Hilbert manifold of space theorem rebuilt with
Mebius space in Gamma Function on Beta function fill into power of rout fundamental
group.
And this result with $AdS_5$ space time in Quantum caos of Minkowsky of manifold in
abel manifold.
Gravity of metric in non-commutative equation is Global differential equation
conbuilt with Kaluza-Klein space. Therefore this mechanism is $T^{\mu\nu}$ tensor
is equal with $R^{\mu\nu}$ tensor. And This moreover inspect with laplace operator
in stimulate with sign of differential operator.
Minus of zone in Add position of manifold is Volume of laplace equation rebuilt with
Gamma function equal with summative of manifold in Global differential equation,
this result with setminus of zone of add summative of manifold, and construct with
locality theorem straight with fundamental group in world line of surface, this
power is boson and fermison of cone in hyper function.

$$ V(\tau) = [f(x),g(x)] \times [f^{-1}(x), h(x)] $$
$$ \Gamma(p,q) = \int e^{-x} x^{1 - t}dx $$
$$ = \beta(p,q) $$
$$ = \pi (f(\chi,x),x) $$
$$ ||ds^2|| = \mathcal{O}(x)[(f(x) \circ g(x))^{\mu\nu}]dx^{\mu}dx^{\nu} $$
$$ = \lim_{x \to \infty} \sum_{k=0}^{\infty}a_k f^k $$
$$ G^{\mu\nu} = {\partial \over \partial f}\int [f(x)^{\mu\nu} \circ G(x)^{\mu\nu}dx^{\mu}dx^{\nu}]^{\mu\nu}dm $$
$$ = g_{\mu\nu}(x)dx^{\mu}dx^{\nu} - f(x)^{\mu\nu}dx^{\mu}dx^{\nu} $$
$$ [i \pi(\chi,x),f(x)] = i \pi f(x) - f(x) \pi(\chi,x) $$
$$ T^{\mu\nu} = (\lim_{x \to \infty} \sum_{k=0}^{\infty}\int \int [V(\tau) \circ S^{\mu\nu}(\chi,x)]dm)
^{\mu\nu}dx^{\mu}dx^{\nu} $$
$$ G^{\mu\nu} = R^{\mu\nu}T^{\mu\nu} $$
$$ \begin{vmatrix} D^m & dx \\
dx & \sigma^m \end{vmatrix} \begin{vmatrix} \cos \theta & - \sin \theta \\
\sin \theta & \cos \theta \end{vmatrix} \begin{vmatrix} x \\
y \end{vmatrix} = \begin{pmatrix} 1 & 0 \\ 
0 & - 1 \end{pmatrix}^{1 \over 2} $$
$$ \sigma^m \begin{bmatrix} \delta(x) & - 1 \\
1 & \epsilon(x) \end{bmatrix}^{1 \over 2} = \begin{pmatrix} i & 0 \\
0 & - i \end{pmatrix} $$
$$ V(M) = {\partial \over \partial f} ({}^N \int [f \smallsetminus M]^{\oplus N})^{\mu\nu} dx^{\mu}dx^{\nu} $$
$$ V(M) = \pi(2 \int \sin^2 dx)\oplus {d \over df}F^M dx_m $$
$$ \lim_{x \to \infty}\sum_{k=0}^{\infty} a_k f^k = \int (F(V)dx_m)^{\mu\nu}dx^{\mu}dx^{\nu} $$
$$ \bigoplus_{k=0}^{\infty}[ f \smallsetminus g] = \vee(M \wedge N) $$
$$ \pi_1 (M) = e^{-f 2 \int \sin^2 x dm} + O(N^{-1}) $$
$$ = [i \pi(\chi,x),f(x)] $$
$$ M \circ f(x) = e^{-f \int \sin x \cos x dx_m} + \log (O(N^{-1})) $$

 Non-Symmetry space time.
 $$ {d \over dL}V(\tau) = {d \over df}\int \int_M {1 \over (x \log x)^2}dx_m + {d \over df}\int \int_M
 {1 \over (y \log y)^{1 \over 2}}dy_m $$
 $$ \epsilon S(\nu) = \square_v \cdot {\partial \over \partial \chi}({}^5 \sqrt{\wedge g^2})d\chi $$
 Differential Volume in $AdS_5$ graviton of fundamental rout of group.
 $$ \wedge (F^m_t)^{''} = {1 \over 12}g_{ij}^2$$
 Quarks of other dimension.
 $$ \pi(V_{\tau}) = {e^{-(\sqrt{\pi \over 16}\log x)}}^{\delta} \times {1 \over (x \log x)} $$
 Universe of rout, Volume in expanding space time.
 $$ {d \over dt}(g_{ij})^2 = {1 \over 24}(F^m_t)^2 $$
 $$ m^2 = 2 \pi T \left({26 - D_n \over 24}\right) $$
 This quarks of mass in relativity theorem, and fourth of universe in three manifold of one dimension surface,
 and also this integrate of dimension in conbuilt of quarks.
 $$ g_{ij} \wedge \pi(\nu_{\tau}) = e^{-2 \pi T |\psi|}[\eta_{\mu\nu} + \bar{h}_{\mu\nu}(x)]dx^{\mu}dx^{\nu} 
 + T^2 d \psi^2 $$
 Out of rout in $AdS_5$ space time.
 $$ ||ds^2|| = g_{ij}\wedge \pi(\nu_v) $$
 $AdS_5$ norm is fourth of universe of power in three manifold out of rout.


  Sphere orbital cube is on the right of hartshorne conjecture by the right equation,
this integrate with fourth of piece on universe series, and this estimate with one of three manifold.
Also this is blackhole on category of symmetry of particles.
These equation conbuilt with planck volume and out of universe is phologram field, this field construct with
electric field and magnitic field stand with weak electric theorem.
This theorem is equals with abel manifold and $AdS_5$ space time.
Moreover this field is anti-brane and brane emerge with gravity and antigravity equation restructed with.
Zeta function is existed with Re pole of ${1 \over 2}$ constance reluctances.
This pole of constance is also existed with singularity of complex fields.
Von Neumann manifold of field is also this means.
$$ ||ds^2|| = \lim_{x \to \infty} [\delta(x)\int \int \int \pi \left( \sum_{k=0}^{\infty}{{}^n\sqrt{p},x \over n}\right)
^{1 \over 2}d \tau]^{\mu\nu} $$
$$ e^{i \theta} = \cos \theta + i \sin \theta $$
$$ e^{x \log x} = x^{{1 \over 2} + iy}, x \log x = \log (\cos \theta + i \sin \theta) $$
$$ = \log \cos \theta + i \log \sin \theta $$
This equation is developed with frobenius theorem activate with logment equation.
This theorem used to
$$ \log(\sin \theta + i \cos \theta)= \log(\sin \theta - i \cos \theta) $$
$$ \log \left({\sin \theta \over i \cos \theta}\right) = -2 R_{ij}, {d \over dt}g_{ij}(t) =  - 2 R_{ij} $$
$$ \mathcal{O}(x) = {\zeta(s) \over \sum\limits_{k=0}^{\infty}a_k f^k} $$
$$ \mathrm{Im} f = \mathrm{ker} f, \chi(x) = {\mathrm{ker}f \over \mathrm{Im}f} $$ $$ H(3) = 2, \nabla H(x) = 2, 
 \pi(x) = 0 $$
 This equation is world line, and this equation is non-integrate with relativity theorem of rout constance.

$$ [f(x)] = \infty, ||ds^2|| = \mathcal{O}(x)[\eta_{\mu\nu} + \bar{h}_{\mu\nu}(x)]dx^{\mu}dx^{\nu} + T^2 d^2 \psi $$
$$ T^2 d^2 \psi = [f(x)] , T^2 d^2 \psi = \lim_{x \to 1}\sum_{k=0}^{\infty}a_k f^k $$

 These equation is equals with M theorem. For this formula with zeta function into one of universe in four dimensions.
Sphere orbital cube is on the surface into $AdS_5$ space time, and this space time is Re pole and Im pole of 
constance reluctances.


  Gauss function is equals with Abel manifold and Seifert manifold. Moreover this function is infinite time,
and this energy is cover with finite of Abel manifold and Other dimension of seifert manifold on the surface.


  Dilaton and fifth dimension of seifert manifold emerge with quarks and Maxwell equation, this power
is from dimension to flow into energy of boson. 
This boson equals with Abel manifold and $AdS_5$ space time, and this space is created with Gauss function.


  Relativity theorem is composited with infinite on D-brane and finite on anti-D-brane, This space time
restructed with element of zeta function.
Between finite and infinite of dimension belong to space time system, estrald of space element have with 
infinite oneselves. Anti-D-brane have infinite themselves, and this comontend with fifth dimension of AdS5 
have for seifert manifold, this asperal manifold is non-move element of anti-D-brane and move element 
of D-brane satisfid with fifth dimension of seifert algebra liner.
Gauss function is own of this element in infinite element.
And also this function is abel manifold.
Infinite is coverd with finite dimension in fifth dimension of AdS5 space time.
Relativity theorem is this system of circustance nature equation.
AdS5 space time is out of time system and this system is belong to infinite mercy.
This hiercyent is endrol of memolite with geniue of element.
Genie have live of telomea endore in gravity accesorlity result.
AdS5 space time out of over this element begin with infinite assentance.
Every element acknowlege is imaginary equation before mass and spiritual envy.

  $$ T^2 d^2 \psi = \lim_{x \to 1}\sum_{k=0}^{\infty}a_k f^k $$
  $$ \lim_{x \to 1}\sum_{k=0}^{\infty}a_k f^k = [T^2 d^2 \psi] $$

  $$ {d \over dL}V(\tau)= {d \over df}\int \int_M({}^5\sqrt{x^2})d\Lambda 
  + {d \over df}\int \int_M {}^N({}^3\sqrt{x})^{\oplus N}d\Lambda $$
  $$ {}^M(\vee(\wedge f \circ g)^N)^{1 \over 2} = {d \over df}\int \int_M{1 \over (x \log x)^2}dx_m
  + {d \over df}\int \int_M{1 \over (y \log y)^{1 \over 2}}dy_m $$
  $$ ||ds^2|| = \mathcal{O}(x)[\eta_{\mu\nu} + \bar{h}_{\mu\nu}(x)]dx^{\mu}dx^{\nu}
  + T^2 d^2 \psi $$
  $$ \mathcal{O}(x) = e^{-2\pi T|\psi|} $$
  $$ G^{\mu\nu} = R_{\mu\nu} T^{\mu\nu} $$
  $$ = - {1 \over 2}\Lambda g_{ij}(x) + T^{\mu\nu} $$

  Fifth dimension of eight differential structure is integrate with one geometry element,
Four of universe also integrate to one universe, and this universe represented with symmetry formula.
Fifth universe also is represented with seifert manifold peices.
All after universe is constructed with six element of circuatance.
This aquire maniculate with quarks of being esperaled belong to.

\
   大域的微分多様体は、世界線を求めている。 
   どの無限階線形微分でも、大域的微分多様体を施すと、世界線になり、
   大域的微分多様体の大域的2重微分多様体は、特異点を求めている。
   トポロジーのホモロジーとコホモロジーがこの性質を担っている。
   $$ e^{{\sqrt{3} \over 2}\log {{\sqrt{3} \over 2}}} $$
   $$ = {\sqrt{3} \over 2}^{{\sqrt{3} \over 2}} $$
   大域的積分多様体は、曲平面を求めている。
   木星の大善が、永遠の生命エネルギーであるのと、地球から宇宙へのロケットの発射の角度が、
   60度より、相似した構造にもなっているのを言えることから、宇宙から異次元への移行の角度も、
   60度と言える。

   $$ \sin \theta = {{\sqrt{3} \over 2}} $$
   は、拡大のグラフも縮小のグラフも、値は、同一である。 
   代数幾何の量子化の因数分解は、加減乗除の式のm,nの
   組み合わせ多様体でのガンマ関数同士の計算になる。
   $$ \left(\bigoplus({i \hbar^{\nabla}})^{\oplus L}+m\right) 
   \left(\bigoplus({i \hbar^{\nabla}})^{\oplus L}+n\right) $$
  $$ = {n^{L+1} \over {n+1}} = t \int (1 - x)^n x^{m - 1} dx $$
  $$ \int x^{m-1}(1 - x)^{n-1} dx $$
  $$ \nabla ({i \hbar^{\nabla}})^{\oplus L}, \bigoplus ({i \hbar^{\nabla}})^{
	   \oplus L}, \square ({i \hbar^{\nabla}})^{\oplus L} $$
   $$ \boxtimes (i \hbar^{\nabla})|_{dx_m}^L, \boxplus (i \hbar^{\nabla}|_
   {dx_m}^L) $$

   宇宙の周りは、マンデルブロー集合でのベータ関数になっている構造である。 




  微分幾何の量子化は、代数幾何の量子化の計算になっている。
  加減乗除が初等幾何であり、大域的微分と積分が、現代数学の代数計算の
  簡易での楽になる計算になっている。
  初等代数の計算は、
  $$ \bigoplus({i \hbar^{\nabla}})^{\oplus L} $$
   $$m \bigoplus({i \hbar^{\nabla}})^{\oplus L} +  
   n\bigoplus({i \hbar^{\nabla}})^{\oplus L} = 
   (m+n) \bigoplus({i \hbar^{\nabla}})^{\oplus L}   $$
   $$ m\bigoplus({i \hbar^{\nabla}})^{\oplus L} -  
   n\bigoplus({i \hbar^{\nabla}})^{\oplus L} = 
   (m-n) \bigoplus({i \hbar^{\nabla}})^{\oplus L}   $$
   $$ \bigoplus({i \hbar^{\nabla}})^{{\oplus L}^m} \times 
      \bigoplus({i \hbar^{\nabla}})^{{\oplus L}^n} =
   \bigoplus({i \hbar^{\nabla}})^{\oplus L^{m+n}} $$
  $$ {{\bigoplus({i \hbar^{\nabla}})^{\oplus L^m}}  \over  
  {\bigoplus({i \hbar^{\nabla}})^{\oplus L^n}}} =  
  \bigoplus({i \hbar^{\nabla}})^{\oplus L^{m\over n}} $$
  大域的計算での微分と積分は、
  $$ \left(\bigoplus({i \hbar^{\nabla}})^{\oplus L}\right)^{df}
  =  {\left(\bigoplus({i \hbar^{\nabla}})^{\oplus L}\right)^{
	  \bigoplus({i \hbar^{\nabla}})^{\oplus L}}}^{'}  $$
  $$ = \bigoplus({i \hbar^{\nabla}})^{\oplus L^{'}} $$
   $$ \int \bigoplus({i \hbar^{\nabla}})^{\oplus L} dx_m $$
   $$ \bigoplus({i \hbar^{\nabla}})^{\oplus L} = n $$
   大域的積分の代数幾何の量子化での計算は、ガンマ関数になっている。
   $$ {n^{L+1} \over {n+1}} = \int e^x x^{1 - t} dx $$
   代数幾何の量子化の因数分解は、加減乗除の式のm,nの
   組み合わせ多様体でのガンマ関数同士の計算になる。
   $$ \left(\bigoplus({i \hbar^{\nabla}})^{\oplus L}+m\right) 
   \left(\bigoplus({i \hbar^{\nabla}})^{\oplus L}+n\right) $$
  $$ = {n^{L+1} \over {n+1}} = t \int (1 - x)^n x^{m - 1} dx $$
  $$ \int x^{m-1}(1 - x)^{n-1} dx $$
  $$ \nabla ({i \hbar^{\nabla}})^{\oplus L}, \bigoplus ({i \hbar^{\nabla}})^{
	   \oplus L}, \square ({i \hbar^{\nabla}})^{\oplus L} $$
   $$ \boxtimes (i \hbar^{\nabla})|_{dx_m}^L, \boxplus (i \hbar^{\nabla}|_
   {dx_m}^L) $$



  $$ x^{{1 \over 2} + iy} = e^{x \log x} $$
  それゆえに、この定数はゼータ関数と微分幾何の量子化を因数にもつ
  素因子分解の式にもなっている。
  Therefore, this product also constructed with differential geometry
  of quantum level equation and zeta function.
  $$ = C $$ 
  そして、この関数はオイラーの定数から広中平祐定理による4重帰納法の
  オイラーの公式からの多様体積分へとのサーストン空間のスペクトラム関数
  ともなっている。
  And this function of Euler product respectrum of focus with
  Heisuke Hironaka manifold in four assembled of integral Euler equation.. 

  $$ C =  b_0 + \cfrac{c_1}{b_1 +
	  \cfrac{c_2}{b_2 +
	  \cfrac{c_3}{b_3 +
	  \cfrac{c_4}{b_4 + \cdots}}}} $$
  この方程式は指数による連分数としての役割も担っている。  
  This equation demanded with continued number of step function.
  $$ x^{{1 \over 2} + iy} = e^{x \log x} $$
  $$ (2.71828)^{2.828} = a^{a + b^{b + c^{c + d^{d \cdots }}}} $$
  $$ = \int e^f \cdot x^{1 - t}dx $$
  This represent is Gamma function in Euler product.
  Therefore this product is zeta function of global differential equation.

  $$ {d \over d f}F(v_{ij},h) = \int {e^{-f}\mathrm{dV}}
  [- \Delta v + \nabla_{i}\nabla_{j}
  v_{ij} - R_{ij}v_{ij} - v_{ij}\nabla_{i}\nabla_{j} + 2 <\nabla f,\nabla h>
  + (R + \nabla f^2)({v \over 2} - h)]$$
  $$ = {d \over d f}F = {2 {\int (R + \nabla_{i}\nabla_{j} f)^2} \over 
  { - (R + \Delta f)}}\mathrm{dm} $$
  これらの方程式は8種類の微分幾何の次元多様体として、そして、これらの
  多様体は曲平面による双対性をも生成している。 
  そして、このガウスの曲平面は、大域的微分多様体と微分幾何の量子化から
  素因数を形成してもいる。
  These equation are eight differential geometry of dimension calyement.
  And these calyement equation excluded into pair of dimension surface.
  This surface of Gauss function are global differential manifold,
  and differential geometry of quantum level.
  $$ F \ge {d \over d f}\int \int{1 \over (x \log x)^2}dx_m 
  + {d \over d f}\int \int{1 \over (y \log y)^{1 \over 2}}dy_m $$

  微分幾何の量子化はオイラーの定数とガンマ関数が指数による連分数
  としての不変性として素因数を形成していて、このガウスの曲平面
  による量子力学における重力場理論は、ダランベルシアンの切断多様体
  がこの大域的切断多様体を付加してもいる。
  Differential geometry of quantum level constructed with Euler product
  and Gamma function being discatastrophed from continued fraction style.
  Gravity equation lend with varint of monotonicity of level expresented
  from gravity of letter varient formula.
  これらの方程式は基本群と大域的微分多様体をエスコートしていもいて、
  ヴェイユ予想がこのダランベルシアンの切断方程式たちから
  輸送のポートにもなっている。ベータ関数とガンマ関数がこれらのフォームラ
  の方程式を放出してもいて、結果、これらの方程式は広中平祐定理の複素
  多様体とグリーシャ教授によるペレルマン多様体からサポートされてもいる。
  この2名の教授は、一つは抽象理論をもう一方は具象理論を説明としている。
  These equation escourted into Global differential manifold and
  fundemental equation.
  Weil's theorem is imported from this equation in gravity of letters.
  Beta function and Gamma function are excluded with these formulas.
  These equation comontius from Heisuke Hironaga of complex manifold
  and Gresha professor of Perelman manifold.
  These two professos are one of abstract theorem and the other of
  visual manifold theorem.


\begin{equation}
 \Delta E = - 2(T- t)|R_{ij} + \nabla_{i} \nabla_{j} f - 
	{1 \over 2(T - t)}g_{ij}|^2
\end{equation}
$\text{Quote group classify equivalent class to own element of group.}$
\begin{equation}
   A = BQ + R 
\end{equation}

\begin{equation}
   [x] = A 
  dx^n = \sum^\infty_{k=0} x^n dx 
\end{equation}

\begin{equation}
   R_{n} = {n! \over (n - r)!} (x^n)' 
\end{equation}

\begin{equation}
   \beta (p,q) = \int^1_0 x^{p-1} (1 - x)^{q-1} 
\end{equation}

\begin{equation}
   Z(T,X) = \mathrm{exp} (\sum^\infty_{m=1}{(q^k T)^m \over m}) 
\end{equation}

\begin{equation}
   Z(T,X) = {P_{1}(T) P_{3}(T) \dots P_{2n-1} \over  
	P_{0}(T) P_{2}(T) P_{4}(T) \dots P_{2n}} 
\end{equation}

\begin{equation}
   |v| = |\int (\pi r^2 + \vec{r})dx|^2 
\end{equation}

\begin{equation}
   \Delta E = \int (\mathrm{div} (\mathrm{rot} E) \cdot e^{-ix \log x}) dx 
\end{equation}


\begin{equation}
   (\nabla \phi)^2 = \int tf(t){df(x) \over e^{-x} t^{x - 1}} dx 
\end{equation}

\begin{equation}
   (\nabla \phi)^2 = \int tf(t)(\Gamma(t)df(x)) dx 
\end{equation}

\begin{equation}
   (\nabla \phi)^2 = {1 \over \Gamma(x + y)} 
\end{equation}

   then these equation decide to class manifold with group.
  Differential group emerge with same element of equation, zero dimension conclude to emerge with all element.
  Constant has with imaginary of number on developed of zeta equation. Weil's Theorem resolved with zeta equation,
these function merge to build in replace equation. Euler constant has with bind of imaginary number and variable 
number.

$\text{	These function is}$

\begin{equation}
	C = \int {1 \over x^s}dx - \log x 
\end{equation}

$\text{Replace equation resolve on zeta function.}$

\begin{equation}
	\int \int {1 \over (x \log x)^2}dx_{m} = {1 \over 2}i 
\end{equation}

\begin{equation}
	\int C dx_m = \int{({\int {1 \over x^s}dx}  + \log x)}d \mathrm{vol}
\end{equation}
\begin{flushright}
$\blacksquare$
\end{flushright}

\begin{equation}
	\int C dx_m = {d \over d f}\int \int {1 \over (y \log y)^{1 \over 2}}
	dy_m + \int \log y d \mathrm{vol}
\end{equation}
\begin{flushright}
$\blacksquare$
\end{flushright}

\begin{equation}
	\int C dx_m = \int \Gamma(\gamma)^{'}dx_m
\end{equation}
\begin{flushright}
$\blacksquare$
\end{flushright}

   These function understood is become of imaginary number,
  which deal with delete line of equation on space of curve.


   Euler number also has with imaginary of constant.
\
ベータ関数をガンマ関数へと渡すと、
   Gamma function sented for Beta function of inspiration with monitonicity
   of component with differential and integral being definitions.
   $$ \beta(p,q) = {{\Gamma(p)\Gamma(q)} \over {\Gamma(p + q)}} $$
   ここで、単体的微分と単体的積分を定義すると
   This area defined with different and integrate of component with
   global topology.
   $$ t = \Gamma(x) $$
   $$ T = \int \Gamma(x) dx $$
   これを使うのに、ベータ関数を単体的微分へと渡すために、
   for this system used for being is conbiniate 
   with beta function for componenet of deprivate eqaution.
   $$ \beta(p,q) = - \int {1 \over {t^2}}dt $$
   大域的微分変数へとするが、
   This point be for global differential variable exchanged,
   $$ T_m = \int \Gamma(x) dx_m $$
   この単体的積分を常微分へとやり直すと、
   This also point be for being retried from ordinary differential computation
   for component of integral manifold.
   $$ T^{'} = {t^{'} \over {{\log t}}}dt + C(\text{Cは積分定数}) $$
   これが、大域的微分多様体の微分変数と証明するためには、
   This exceed of proof being for being defined with deprivate variable of
   global differential formula, this successed for true is,
   $$ dx_m = {1 \over {\log x}}dx, dx_m = {(\log x^{-1})}^{'}  $$
   これより、単体的微分が大域的部分積分へと置換できて、
   Therefore, this exchanged from mononotocity deprivation from global
   parital integral manifold is able to, 
   $$ T^{'} = \int \Gamma(\gamma)^{'} dx_m $$

   微分幾何へと大域的積分と大域的微分多様体が、加群分解の共変微分へと
   書き換えられる。
   After all, these exchanged of monotonicity deprivation successed from
   being catastrophe of summativate of partial and assemble of deprivations
   for differential geometory of quantum level to global integral and
   differential manifold.
   $$ \bigoplus T^{\nabla}  = \int T dx_m, \delta(t) = t dx_m  $$


   this exluded of being conclution are which beta function
   evaluate with mononicity from ordinary differential equation
   be resulted from component of deprivation and integral expalanations.
   This cirtutation be resenbled to define with global topoloty of
   extention extern of deprivate and integral of manifolds
   estourced with quantum level of differential geometry be proofed with
   all of equation anbrabed from Euler product.
   これらをまとめると、ベータ関数を単体的関数として、
   これを常微分へと解を導くと、単体的微分と単体的積分
   へと定義できて、これより、大域的微分多様体と大域的積分多様体
   へとこの関数が存在していることを証明できるのを上の式たちから
   わかる。この単体的微分と単体的積分から微分幾何が書けることがわかる。



   虚数の虚数倍した値が超越数のx倍と同じとすると、
   超越数の$2 \pi m$倍が$n=1$でiとなると定義すると、
   次の式たちが導かれる。　
   Imaginary pole circlate with twigled of pole in step function
   from Naipia number of assembled from equalation of defined are
   escourted to be defined with next equations.


   These defined equation are climbate with idea of equation from
   Caltan of imaginary number of circulation.
   カルタンを超えているアイデアと数式たちでもある。

   $$ i^i = (\sqrt{i})^{\sqrt{i}} = e^{x \log x} $$
   $$ e^x = i, e^{2\pi m} = i^n, {d \over d x}e^{2\pi m} = i^n $$
   $$ e^{i \theta} = \cos \theta + i \sin \theta, e^{2 \pi m} = i^n $$
   $$ 2 \pi m = n, e = i \text{($e=i$となる$e^{i \theta} = \cos \theta 
   + i \sin \theta $
   の範囲で$2\pi m$ がある。)}$$
   $$ {d \over d i}i = i^i = e^{x \log x}, m={1 \over 2 \pi},l=2\pi r, 
   \pi = {l \over 2r} $$
   $$ ||ds^2|| = e^{-2 \pi T||\psi||}[\eta_{\mu\nu} + \bar{h}(x)]
   dx^{\mu}dx^{\nu} + T^2 d^2 \psi $$
   $$ e^x = i, e^{2\pi m} = i^n, e = i $$
   $$ {{[\eta_{\mu\nu} + \bar{h}(x)]dx^{\mu}dx^{\nu}}/i} = H \Psi
   = i \hbar \psi, H \Psi = {1 \over i}[H, \Psi] $$
   ハイゼンベルク方程式が$ AdS_5 $多様体の原子レベルの方程式も表せられる。 
   微分幾何の量子化の式は、
   Hisenbelg equation are represeted from AdS5 manifold to particle level
   equalation of quantum levle of differential geometry entranced.
   $$ \bigoplus {({i \hbar} ^{\nabla})}^{\oplus L} = e^{{i x \log x}^{
	   {e^{i x \log x}}^{'}}} $$
    $$ H\Psi = \bigoplus {H \Psi \over \nabla L} $$

   この式が、アカシックレコードの子どもの式である。２種類の表し方である。
   計算方法が、大域的微分多様体の求め方と同じである。

   移項して、
   それを展開していき、
    これらより、まとめた式が、
    $$ n!^{-1} = \bigoplus \nabla L $$
    これは、積分再発見法と同じであり、
    $$ n!^{-1}  = \int T^{{\mu\nu}^{'}} dV - \int T^{\mu\nu} dV $$
    これもカルーツァ・クライン空間と同じ式であり、

    $$ = \int f{(x)^{-1}} xf(x) - f(x) $$
    これに条件をつけると、
    $$ \exists x = a, 1 - \exists \int x dx, \Lambda = \int A dx $$
    これより、次の式になり、
    $$ \nabla f(x) = {1 \over n!} \int \int \cdots \int dx $$

    $$ \sum \int (\int \cdots A\int dz)/n! $$
    ウィッテン方程式により、
    $$ M^2 a^{(m)} = k a^{(m)} $$
    $$ ZM^3 = \int dAe^{ki/4\pi}\mathcal{L}M^3 $$
    $$ Z_M = \langle M_1 | M_2 \rangle $$
    ゼータ関数へといける。
    $$ e^{x} = \sum {x^n \over n!} (r = \infty) $$
    これらは、大域的微分多様体の部品たちになっている。
    $$ {d \over df}\int dx_m = {1 \over n!}\sum x^n $$


   Euler productは、
   $$ \int^{n+1}_{n}{dx \over x} < {1 \over n}, (n \ge 1) $$
   $$ 1 + {1 \over 2} + \cdots + {1 \over n} > \int^{n+1}_{1}{dx \over x}
   = \log(n+1), $$
   $$ 1 + {1 \over 2} + \cdots + {1 \over n} - \log x > \log{{(n+1)} 
   \over n} > 0, $$
   $$ {1 \over {n+1}} < \int^{n+1}_{0} {dx \over x} = \log(n+1) - \log n. $$
   $$ \lim_{n \to \infty} \left(1 + {1 \over 2} + \cdots + {1 \over n}
   - \log n \right) = C $$
   $$ = \int (\int {1 \over x^s}dx - \log x)\mathrm{dvol} = \int C dx_m $$
   Cの値は、$\text{C = 0.5772156$\cdots$} $である。

  オイラーの定数は、次の式で成り立っている。
  $$ C = \int{1 \over x^s}dx - \log x $$
  その式を単体積分をすると、
  $$ = \int {\left({\int {1 \over x^s}}dx - \log x\right)}\mathrm{dvol} $$
  ゼータ関数を多重積分すると、　
  $$ = \int \int {1 \over x^s}dx - \int {\log x} \mathrm{dvol} $$
  ここで、対数方程式は、多様体積分がガンマ関数を微分した方程式と
  同値より、
  $$ F = \Gamma = \int e^{-x}x^{1-t}dx $$
  $$ = {d \over df}\int \int {1 \over (y \log y)^{1 \over 2}}dx_m
  - \int e^{-x}x^{1 - t} \log x dx_m $$
  $$ \int x^{1 - t}e^{-x}dV = \int x^{1 - t}dm $$
  $$ \int x^{1 - t}e^{-x}dV = \int x^{1 - t}\mathrm{dvol} $$
  $$ f = \gamma = \Gamma^{'} = \int e^{-x}x^{1 - t} \log x dx $$
  $$ = {d \over d \gamma}\Gamma^{-1} - 
  \int e^{-x}x^{1 - t}\log x dx_m $$
  $$ = {d \over d \gamma}\Gamma^{-1} - \left(\gamma\right)^{\gamma^{'}} $$
  これらより、大域的微分方程式のヒッグス場方程式の逆三角関数の双曲多様体
  に対極する、双曲多様体の余弦定理と同値になる。
  $$ = e^{-f} - e^{f} $$
  $$ = 2 \cos (i x \log x) $$
  結局は、微分幾何の量子化がガンマ関数でもあり、その逆関数はゼータ関数
  でもあり、そして、オイラーの定数を多様体微分をすると、ヒッグス場の
  方程式は正弦定理の逆三角関数であり、このヒッグス場の
  方程式の逆三角関数が、余弦定理の
  双曲多様体として、オイラーの定数になる。
  オイラーの定数は、このガンマ関数から、無理数と証明される。

    カルーツァ・クライン空間の方程式は、
    $$ ds^2 = g_{\mu\nu}(x)dx^{\mu}dx^{\nu} + \psi^2(x)(dx^2 + \kappa^2
    A_{\mu}(x)dx^{\nu})^2 $$
    この式は、
    $$ ds^2 = -N(r)^2dt^2 + \psi^2(x)(dr^2 + r^2d\theta^2) $$
    $$ ds^2 = -dt^2 + r^{-8\pi Gm}(dr^2 + r^2d\theta^2) $$
    とまとまり、
    $$ dx^2 = g_{\mu\nu}(x)(g_{\mu\nu}(x)dx^2 - dxg_{\mu\nu}(x)) $$
    $$ dx = (g_{\mu\nu}(x)^2dx^2 - g_{\mu\nu}(x)dx g_{\mu\nu}(x))^{1 \over 2} $$
    と反重力と正規部分群の経路和となり、　
    $$ \pi(\chi, x) = i \pi (\chi, x)f(x) - f(x)  \pi (\chi, x) $$
    と基本群にまとまる。
    シュワルツシルト半径は、
    $$ ds^2 = -(1 - {r_s \over r})c^2 dt^2 + {1 \over {{1 - {r_s \over r}}}} 
    dr^2 + r^2 d \theta^2 + r^2 \sin^2 \theta d \psi^2 $$
    これは、カルーツァ・クライン空間と同型となる。
    一般相対性理論は、
    $$ R^{\mu\nu} + {1 \over 2}g_{ij}\Lambda = \kappa T^{\mu\nu} $$
    この多様体積分は、
    $$ \int \kappa T^{\mu\nu}\mathrm{dvol} = \int  \left(R^{\mu\nu} + 
    {1 \over 2}g_{ij}\Lambda\right)\mathrm{dvol}  $$
    ガンマ関数のおける大域的微分多様体は、これと同型により、
    $$ \int \Gamma \cdot {d \over {d \gamma}}{\Gamma dx_m} \le \int \Gamma
    dx_m + {d \over d \gamma}\Gamma $$
    $$ = \int \Gamma(\gamma)^{'}dx_m $$
    オイラーの定数の多様体積分は、
    $$ \int \left({{\int {1 \over x^s}}dx - \log x} \right)\mathrm{dvol} $$
    この解は、シュワルツシルト半径と同型より、
    $$ = e^f - e^{-f} \le {e^f + e^{-f}} $$
    次元の単位は多様体より、大域的微分多様体とオイラーの定数
    の多様体積分の加群分解は、オイラーの公式の三角関数の虚数の度解より、
    $$ {d \over df}F + \int C dx_m = 2 (\cos (i x \log x) 
    - i \sin (i x \log x)) $$
    すべては、大域的微分多様体の重力と反重力方程式に行き着き、
    $$ {d \over df}F \ge {d \over df}\int \int{1 \over (x \log x)^2}dx_m
    + {d \over df}\int \int{1 \over (y \log y)^{1 \over 2}}dy_m $$
    と求まる。

    これが、重力と反重力単独では、磁気単極子で、ド・ブロイのアイデアであり、
    両方では、磁気双極子で、南部・ゴールドストーンボソン粒子である。
    ヒッグス場の式でもある。

    北半球の磁気と南半球に磁気一セットで一種である。

    これらは、ハイゼンベルクとシュレーディンガー方程式のディラック作用素が、
    微分幾何からの作用素関数から、同一と言える。

    
    一般相対性理論における多様体積分とオイラーの定数の多様体積分
    が同型と言えることを上の式たちは述べている。

    このJones多項式が、金融市場の相場価格を決めるオプション方程式であり、
    流体力学による熱エントロピー値の熱エネルギー量が、
    いろんな機材に使われている。
    量子コンピュータにおける論理素子の回路の設計にも使われている。
    人工知能の論理素子と因数分解における量子暗号にも使われている。

    この一般相対性理論における多様体積分がJones多項式となり、
    金融市場で使われることを読んだのが、イスラエル国家らしい。

    $$ H\Psi = \bigoplus {\left(i \hbar^{\nabla} \right)}^{\oplus L} $$
    $$ = {1 \over 2}i e^{i \Hat{H}} $$
    $$ (i \hbar)^{'} = (- e^{i \Hat{H}})^{'} $$
    $$ = - i e^{i \Hat{H}} $$
    $$ \psi(x) = e^{-i \Hat{H}t}, \bigoplus (i \hbar ^{\nabla})^{\oplus L} 
    = {{{1 \over 2} e^{i \Hat{H}}}^{(-i e^{i \Hat{H}})}} $$
    $$ = ({1 \over 2}f)^{-i f}$$
    $$ = ({1 \over 2})^{-i f} \cdot e^{-x \log x} \cdot (f)^{i} $$
    $$ = \int e^{- x} x^{t - 1}dx, {d \over d \gamma}\Gamma = e^{-x \log x}  $$
     
    $$ f = x, i = t, {1 \over 2} = a, \bigoplus a^{-t x}x^t [I_m]
     \cong \int e^{-x} x^{t - 1}dx $$
     $$ |\psi(t)\rangle_s = e^{-i\Hat{H} t}|\Psi\rangle_{H},
    \Hat{A}_s = \Hat{A}_H(0) $$
    $$ |\Psi(t)\rangle_s \to {d \over dt} $$
    $$ i {d \over dt}|\psi (t) \rangle_s = \Hat{H}|\psi(t)\rangle_s $$
    $$ \langle\Hat{A}(t)\rangle = \langle\Psi(t)|\Hat{A}(0)|\Psi(t)\rangle  $$
    $$ {d \over dt}\Hat{A} = {1 \over i}[\Hat{A},H] $$
    $$ \Hat{A}(t) = e^{i\Hat{H}t}\Hat{A}(0)e^{-i\Hat{H}t} $$
    $$ \lim_{\theta \to 0} \begin{pmatrix} \sin \theta \\
    \cos \theta \end{pmatrix} \begin{pmatrix} \theta & 1 \\
    1 & \theta \end{pmatrix} \begin{pmatrix} \cos \theta \\
    \sin \theta \end{pmatrix} = \begin{pmatrix} 1 & 0 \\
    0 & - 1 \end{pmatrix} $$
    $$ f^{-1}(x) x f(x) = I^{'}_m, I^{'}_m = [1,0] \times [0,1] $$

    $$ {x + y} \ge \sqrt{x y} $$
    $$ {x^{{1 \over 2} + iy} \over e^{x \log x}} = 1 $$
    $$ \mathcal{O}(x) = \nabla_i \nabla_j \int e^{{2 \over m}\sin \theta 
    \cos \theta } \times {{N \mathrm{mod}(e^{x \log x})} \over {O(x)(x +
    \Delta |f|^2)^{1 \over 2}}} $$
    $$ x \Gamma(x) = 2 \int|\sin 2 \theta|^2 d \theta, \mathcal{O}(x)
    = m(x)[D^2 \psi] $$
    $$ i^2 = (0,1) \cdot (0,1), |a||b|\cos \theta = - 1 $$
    $$ E = \mathrm{div}(E, E_1) $$
    $$ \left({\{f,g\}} \over [f,g] \right) = i^2, E = mc^2, I^{'} = i^2 $$

    この式たちは、微分幾何の量子化から計算方法としての形式の
    大域的微分多様体の大域的部分積分が,大域的偏微分方程式としての式の
    形からわかるのも、微分幾何の量子化の仕組みとしての計算式から
    導けられるのも、すべては未来の私の子の情報から読み取れたと
    言える。
 
    $$ H \Psi = \bigoplus (i \hbar ^{\nabla})^{\oplus L} $$
    $$ = \bigoplus {{ H \Psi} \over \nabla L} $$
    $$ = e^{x \log x} = x^{(x)^{'}} $$
    however,
    $$ {d \over df}F = m(x) $$
    and gravity equation in being abled to existed with quantum physics,
    and this physics is also represented with quantum level of differential
    geometry.
    $$ {d \over dt}\psi(t) = \hbar $$
    $$ = {1 \over 2}i e^{i \Hat{H}} $$
    $$ (i \hbar)^{'} = (- e^{i \Hat{H}})^{'} $$
    $$ = - i e^{i \Hat{H}} $$
    $$ \psi(x) = e^{-i \Hat{H}t}, \bigoplus (i \hbar ^{\nabla})^{\oplus L} 
    = {{{1 \over 2} e^{i \Hat{H}}}^{(-i e^{i \Hat{H}})}} $$
    $$ = ({1 \over 2}f)^{-i f}$$
    $$ = ({1 \over 2})^{-i f} \cdot e^{-x \log x} \cdot (f)^{i} $$
    $$ = \int e^{- x} x^{t - 1}dx, {d \over d \gamma}\Gamma = e^{-x \log x}  $$
     
    $$ f = x, i = t, {1 \over 2} = a, \bigoplus a^{-t x}x^t [I_m]
     \cong \int e^{-x} x^{t - 1}dx $$
     Quantum level of differential geomerty is also resulted with
    Gamma function of global differential manifold of values, Heisenberg
    equation is alos resulted with global integrate manifold and this
    quantum level of differential geometry is manifold of deprive of formula.
    $$ |\psi(t)\rangle_s = e^{-i\Hat{H} t}|\Psi\rangle_{H},
    \Hat{A}_s = \Hat{A}_H(0) $$
    $$ |\Psi(t)\rangle_s \to {d \over dt} $$
    $$ i {d \over dt}|\psi (t) \rangle_s = \Hat{H}|\psi(t)\rangle_s $$
    $$ \langle\Hat{A}(t)\rangle = \langle\Psi(t)|\Hat{A}(0)|\Psi(t)\rangle  $$
    $$ {d \over dt}\Hat{A} = {1 \over i}[\Hat{A},H] $$
    $$ \Hat{A}(t) = e^{i\Hat{H}t}\Hat{A}(0)e^{-i\Hat{H}t} $$
    $$ \lim_{\theta \to 0} \begin{pmatrix} \sin \theta \\
    \cos \theta \end{pmatrix} \begin{pmatrix} \theta & 1 \\
    1 & \theta \end{pmatrix} \begin{pmatrix} \cos \theta \\
    \sin \theta \end{pmatrix} = \begin{pmatrix} 1 & 0 \\
    0 & - 1 \end{pmatrix} $$
    $$ f^{-1}(x) x f(x) = I^{'}_m, I^{'}_m = [1,0] \times [0,1] $$

    $$ {x + y} \ge \sqrt{x y} $$
    $$ {x^{{1 \over 2} + iy} \over e^{x \log x}} = 1 $$
    $$ \mathcal{O}(x) = \nabla_i \nabla_j \int e^{{2 \over m}\sin \theta 
    \cos \theta } \times {{N \mathrm{mod}(e^{x \log x})} \over {O(x)(x +
    \Delta |f|^2)^{1 \over 2}}} $$
    $$ x \Gamma(x) = 2 \int|\sin 2 \theta|^2 d \theta, \mathcal{O}(x)
    = m(x)[D^2 \psi] $$
    $$ i^2 = (0,1) \cdot (0,1), |a||b|\cos \theta = - 1 $$
    $$ E = \mathrm{div}(E, E_1) $$
    $$ \left({\{f,g\}} \over [f,g] \right) = i^2, E = mc^2, I^{'} = i^2 $$
     Gamma function of global differential equation is first of differential
    function deprivate with ordinary differential equations, more also
    universe and other dimensiion of Higgs field emerge with zeta function,
    more spectrum focus is three dimension of manifold of singularity or pair
    theorem, and universe and other dimension of each value equals.
    This reverse of circle function of manifolds is also Higgs fields of 
    dense of entropy and result equation. And differential structure of quantum
    level in gravity is poled with zeta function and quantum group in high 
    result values. Zeta funtion and quantum group is Global differential
    manifold of result in values.
    $$ {d \over d \gamma} \Gamma = m(x) $$
    $$ = e^{- x \log x} $$
    $$ \sin i x = {{e^{-x} + e^{x}} \over 2i} $$
    $$ {d \over df}F = m(x) = e^{f} + e^{-f} $$
    $$ = 2 i \sin (i {x \log x}) $$


   Beta function is,
      $$ \beta(p,q) = \int x^{1 - t} (1 - x)^{t}dx = \int t^{x}(1 - t)^{x - 1}
      dt $$
      $$ 0 \le y \le 1, \int^{1}_{0} x^{10}(1 - x)^{20}dx = B(11,21) $$
      $$ = {{\Gamma(11)\Gamma(21)} \over \Gamma(32)} = {{10! 20!} \over {
	      31!}} = {1 \over {931395465}} $$
       $$ {1 \over {931395465}} \cong {1 \over 9} = {1 \over {1 - x}} $$
	 $$     = {1 \over {1 - z}} =
	      \sum_{k=0}^{\infty}z^k = {1 \over {1 + z^2}} = \sum_{
		      k=0}^{\infty} (-1)^k z^{2k} $$
      $$ f(x) = \sum_{k=0}^{\infty} a_k z^k $$
      $$ {{d^n y} \over dx^n} = n! y^{n+1} $$
      $$ f^{(0)}(0) = n!f(0)^{n+1} = n! $$
      $$ f(x) \cong \sum_{k=0}^{\infty}x^n = {1 \over {1 - x}} $$
      In example script is,
      $$ {{dy} \over dx} = y^2, {1 \over y^2}\cdot {dy \over dx} = 1 $$
      $$ \int {1 \over y^2}{dy \over dx}dx = \int{dy \over y^2} = - {1 \over y
      } $$
      $$ - {1 \over y} = x - C, y = {1 \over {C - x}} $$
      $$ \exists x = 0, y = 1 $$
      $$\text{in first value condition compute with} $$
      $$ \text{result, C consumer sartified,} $$
      $$ y = {1 \over {1 - x}} $$
      This value result is concluded with native function from Abel manifold.

      Abelian enterstane with native funtion meltrage from Daranvelcian
      equation on Beta function.
 

   $$ H\Psi = \bigoplus {H \Psi \over \nabla L} $$
  　
 $$ \Gamma = \int e^{-x}x^{1 - t}dx $$
   $$ \gamma = \int e^{-x}x^{1 - t}\log x dx $$
   $$ = \int \Gamma(\gamma)^{'} dx_m$$
   $$ = {d \over d \gamma}\Gamma $$

   Eight differential geometry are each intersect with own level of
   concept from expalanation of Euler product system.
   This component of three manfold sergeried with geometry of destroy and
   desect with time element of Stokes equation.
   ８種類の微分幾何では、それぞれの時間の固有値が違う。
   閉3次元多様体上で微分幾何の切除、分解によって時間の性質が決まる。
   This manifold gut theory from described with zeta function to
   catastrophe for non tree of routs result on sergery of space system.
   閉３次元多様体に統合されると、ゼータ関数になり、各幾何に分解された場合
   に、非分岐から、この上、sergeryの結果が決まる。

   各微分構造においての時間発展での熱エネルギーの変化は$E=mc^2, E=m^2 c^2 $
   のthree manifoldが微分幾何構造の変化にそれぞれ対応する。
   Each geometry point interacte with exchange of three manifold in
   Seifert structure from special relativity of equation for
   heat energy fluentations on time developyed from this extermaination.

   $$ \nabla ({i \hbar^{\nabla}})^{\oplus L}, \bigoplus ({i \hbar^{\nabla}})^{
	   \oplus L}, \square ({i \hbar^{\nabla}})^{\oplus L} $$
   These operator of equation on summatative manifolds from emerged with
   element of particle conclution, Euler product is resulted from this
   operator expalanation of locality insectations.
   これらの作用素が加減乗除の生成式の元、Euler Productの結論による
   作用素生成の論理素子でもある。
   Moreover, this eight geometry of differential operator are
   constructed with four pattern of Jones manifold from
   summatative formula and this system extate with special relativity
   references. This decieved of elemet on summatative equation routed
   to internext in real and imaginary pole on complex dimension,
p  and this dimension explation with Stokes theorem defined too.
   And this defined circutation of Yacobi matrix is represected with
   Knot theory with anstate with Jones manifold.
   その上に、8種類の微分幾何は、Jones多項式の4パターンでの構成される
   差分と加群方程式から、特殊相対性理論をも思わせる、この差分方程式
   $$ {d \over d \gamma}\Gamma = e^{f} + e^{-f} \ge e^{f} - e^{f} $$
   から、8種類の代数幾何が、ストークスの定理と同じく、このJones多項式
   の固有周期が、すべての時間の流れの基軸をも内包してもいることが、
   わかる。このJones多項式は、結び目の理論をも、この周期が言い表してもいる。
   $$ {d \over d \gamma}\Gamma + \int C dx_m = 2(\cos(i x \log x) - i \sin
   (i x \log x)) $$
    $$ e^{- \theta} = \cos(i x \log x) - i \sin(i x \log x)$$
    $$ \int \Gamma(\gamma)^{'}dx_m = 2(\cos (i x \log x) 
    - i \sin(i x \log x)) $$

 カルーツァ・クライン空間の方程式は、
    $$ ds^2 = g_{\mu\nu}(x)dx^{\mu}dx^{\nu} + \psi^2(x)(dx^2 + \kappa^2
    A_{\mu}(x)dx^{\nu})^2 $$
    この式は、
    $$ ds^2 = -N(r)^2dt^2 + \psi^2(x)(dr^2 + r^2d\theta^2) $$
    $$ ds^2 = -dt^2 + r^{-8\pi Gm}(dr^2 + r^2d\theta^2) $$
    とまとまり、
    $$ dx^2 = g_{\mu\nu}(x)(g_{\mu\nu}(x)dx^2 - dxg_{\mu\nu}(x)) $$
    $$ dx = (g_{\mu\nu}(x)^2dx^2 - g_{\mu\nu}(x)dx g_{\mu\nu}(x))^{1 \over 2} $$
    と反重力と正規部分群の経路和となり、　
    $$ \pi(\chi, x) = i \pi (\chi, x)f(x) - f(x)  \pi (\chi, x) $$
    と基本群にまとまる。
    シュワルツシルト半径は、
    $$ ds^2 = -(1 - {r_s \over r})c^2 dt^2 + {1 \over {{1 - {r_s \over r}}}} 
    dr^2 + r^2 d \theta^2 + r^2 \sin^2 \theta d \psi^2 $$
    これは、カルーツァ・クライン空間と同型となる。
    一般相対性理論は、
    $$ R^{\mu\nu} + {1 \over 2}g_{ij}\Lambda = \kappa T^{\mu\nu} $$
    この多様体積分は、
    $$ \int \kappa T^{\mu\nu}\mathrm{dvol} = \int  \left(R^{\mu\nu} + 
    {1 \over 2}g_{ij}\Lambda\right)\mathrm{dvol}  $$
    ガンマ関数のおける大域的微分多様体は、これと同型により、
    $$ \int \Gamma \cdot {d \over {d \gamma}}{\Gamma dx_m} \le \int \Gamma
    dx_m + {d \over d \gamma}\Gamma $$
    $$ = \int \Gamma(\gamma)^{'}dx_m $$
    オイラーの定数の多様体積分は、
    $$ \int \left({{\int {1 \over x^s}}dx - \log x} \right)\mathrm{dvol} $$
    この解は、シュワルツシルト半径と同型より、
    $$ = e^f - e^{-f} \le {e^f + e^{-f}} $$
    次元の単位は多様体より、大域的微分多様体とオイラーの定数
    の多様体積分の加群分解は、オイラーの公式の三角関数の虚数の度解より、
    $$ {d \over df}F + \int C dx_m = 2 (\cos (i x \log x) 
    - i \sin (i x \log x)) $$
    すべては、大域的微分多様体の重力と反重力方程式に行き着き、
    $$ {d \over df}F \ge {d \over df}\int \int{1 \over (x \log x)^2}dx_m
    + {d \over df}\int \int{1 \over (y \log y)^{1 \over 2}}dy_m $$
    と求まる。


  ８種類の微分幾何での、それぞれの時間の固有値の動静は、
  ストークスの流体力学での、熱エネルギーの各流れの仕方で表されている。
  種数１の2種が種数０の球体に両辺で交わっていると、
一般相対性理論における重力理論は大域的微分多様体と積分多様体についての 単体量を求めるための空間の対数関数における不変性を記述するプランクスケール と異次元の宇宙におけるスケール、ウィークスケールと言われているanti-D-brane として、この２種類の計量をガンマ関数とベータ関数としてオイラーの定数 とそれによる連分数が微分幾何の量子化と数式の値を求めると同じという 予知と推測値から、広中平祐定理の４重帰納法のオイラーの公式による 多様体積分と同類値として、ペレルマン多様体がサーストン空間に成り立つ と同じく、広中平祐定理がヒルベルト空間で成り立つとしての、 これら２種の多様体の場の理論が、ゼータ関数がAdS5多様体で成り立つことと、 不変量として、ゼータ関数がこの3種の多様体の場の理論のバランスをとる
  理論として言えることである。
$$ ||ds^2|| = \lim_{x \to \infty} [\delta(x) \int \int \int \pi \left( \sum_{k=0}^{\infty}
		{{}^n\sqrt{p},x \over n} \right)
^{1 \over 2}d \tau]^{\mu\nu} $$
  This norm is component of fields in Hilbert manifold of space theorem rebuilt with
Mebius space in Gamma Function on Beta function fill into power of rout fundamental
group.

  　 Gravity of general relativity theory describe with Cutting of space
in being discatastrophed from Global differential and integral manifold
of scaled levetivity in plank scal and the other vector of universe
scale, these two scale inspectivity of Gamma and Beta function escourted
with these manifold experted in result of Differentail geometry of
quantum level manifold equal with Euler product and continue parameter,
moreover this product is relativity of Hironaka theorem in Four assembled
integral of Euler equation.

 重力場理論の式は、
  Gravity equation is
  $$ \square = - {16 \pi G \over c^4}T^{\mu\nu} $$
  This equation quated with being logment of formula, and
  this formula divided with universe of number in prime zone,
  therefore, this dicided with varint equation is monotonicity of
  being composited with Weil's theorem united for Gamma function.
  この式は、対数を宇宙における数により求める素数分布論として、この
  大域的積分分断多様体がガンマ関数をヴェイユ予想を根幹とする単体量
  として決まることに起因する商代数として導かれる。
  $$ \scalebox{1}[2]{\alearletter} = 
  {{{{8 \pi G} \over {c^4}}T^{\mu\nu}}/ \log x} $$
  $$ {}^{t}\scalebox{1}[2]{\variiint}\mathrm{cohom D_{\chi}[I_m]}$$
  $$ = \oint {(px^n + qx + r)^{\nabla l}} $$
  $$ {d \over d l}L(x,y) = 2 \int ||\sin 2x||^{2}d \tau $$
  $$ {d \over d \gamma}\Gamma $$
  この関数は大域的微分多様体としてのアカシックレコードの合流地点として、
  タプルスペースを形成している。
  This function esterminate with acasic record of global differential
  manifolds.
  $$ = [i \pi (\chi, x), f(x)] $$
  それにより、この多様体は基本群をアカシックレコードの相対性としての
  存在論の実存主義から統合される多様体自身としてのタプルスペースの
  池になっている。
  And this manifold from fundemental group esterminate with
  also this manifold estimate relativity of acasic record.
  $$ {d \over d \gamma}\Gamma = \int C dx_m = \int(\int {1 \over x^s}dx
  - \log x)\mathrm{dvol} $$
  また、このアカシックレコードはオイラーの定数のラムダドライバー
  にもなっている。
  More also this record tupled with lake of Euler product.

  微分幾何の量子化は、代数幾何の量子化の計算になっている。
  加減乗除が初等幾何であり、大域的微分と積分が、現代数学の代数計算の
  簡易での楽になる計算になっている。
  初等代数の計算は、
  $$ \bigoplus({i \hbar^{\nabla}})^{\oplus L} $$
   $$m \bigoplus({i \hbar^{\nabla}})^{\oplus L} +  
   n\bigoplus({i \hbar^{\nabla}})^{\oplus L} = 
   (m+n) \bigoplus({i \hbar^{\nabla}})^{\oplus L}   $$
   $$ m\bigoplus({i \hbar^{\nabla}})^{\oplus L} -  
   n\bigoplus({i \hbar^{\nabla}})^{\oplus L} = 
   (m-n) \bigoplus({i \hbar^{\nabla}})^{\oplus L}   $$
   $$ \bigoplus({i \hbar^{\nabla}})^{{\oplus L}^m} \times 
      \bigoplus({i \hbar^{\nabla}})^{{\oplus L}^n} =
   \bigoplus({i \hbar^{\nabla}})^{\oplus L^{m+n}} $$
  $$ {{\bigoplus({i \hbar^{\nabla}})^{\oplus L^m}}  \over  
  {\bigoplus({i \hbar^{\nabla}})^{\oplus L^n}}} =  
  \bigoplus({i \hbar^{\nabla}})^{\oplus L^{m\over n}} $$
  大域的計算での微分と積分は、
  $$ \left(\bigoplus({i \hbar^{\nabla}})^{\oplus L}\right)^{df}
  =  {\left(\bigoplus({i \hbar^{\nabla}})^{\oplus L}\right)^{
	  \bigoplus({i \hbar^{\nabla}})^{\oplus L}}}^{'}  $$
  $$ = \bigoplus({i \hbar^{\nabla}})^{\oplus L^{'}} $$
   $$ \int \bigoplus({i \hbar^{\nabla}})^{\oplus L} dx_m $$
   $$ \bigoplus({i \hbar^{\nabla}})^{\oplus L} = n $$
   大域的積分の代数幾何の量子化での計算は、ガンマ関数になっている。
   $$ {n^{L+1} \over {n+1}} = \int e^x x^{1 - t} dx $$
   代数幾何の量子化の因数分解は、加減乗除の式のm,nの
   組み合わせ多様体でのガンマ関数同士の計算になる。
   $$ \left(\bigoplus({i \hbar^{\nabla}})^{\oplus L}+m\right) 
   \left(\bigoplus({i \hbar^{\nabla}})^{\oplus L}+n\right) $$
  $$ = {n^{L+1} \over {n+1}} = t \int (1 - x)^n x^{m - 1} dx $$
  $$ \int x^{m-1}(1 - x)^{n-1} dx $$
  $$ \nabla ({i \hbar^{\nabla}})^{\oplus L}, \bigoplus ({i \hbar^{\nabla}})^{
	   \oplus L}, \square ({i \hbar^{\nabla}})^{\oplus L} $$
   $$ \boxtimes (i \hbar^{\nabla})|_{dx_m}^L, \boxplus (i \hbar^{\nabla}|_
   {dx_m}^L) $$



  $$ x^{{1 \over 2} + iy} = e^{x \log x} $$
  それゆえに、この定数はゼータ関数と微分幾何の量子化を因数にもつ
  素因子分解の式にもなっている。
  Therefore, this product also constructed with differential geometry
  of quantum level equation and zeta function.
  $$ = C $$ 
  そして、この関数はオイラーの定数から広中平祐定理による4重帰納法の
  オイラーの公式からの多様体積分へとのサーストン空間のスペクトラム関数
  ともなっている。
  And this function of Euler product respectrum of focus with
  Heisuke Hironaka manifold in four assembled of integral Euler equation.. 

  $$ C =  b_0 + \cfrac{c_1}{b_1 +
	  \cfrac{c_2}{b_2 +
	  \cfrac{c_3}{b_3 +
	  \cfrac{c_4}{b_4 + \cdots}}}} $$
  この方程式は指数による連分数としての役割も担っている。  
  This equation demanded with continued number of step function.
  $$ x^{{1 \over 2} + iy} = e^{x \log x} $$
  $$ (2.71828)^{2.828} = a^{a + b^{b + c^{c + d^{d \cdots }}}} $$
  $$ = \int e^f \cdot x^{1 - t}dx $$
  This represent is Gamma function in Euler product.
  Therefore this product is zeta function of global differential equation.

  $$ {d \over d f}F(v_{ij},h) = \int {e^{-f}\mathrm{dV}}
  [- \Delta v + \nabla_{i}\nabla_{j}
  v_{ij} - R_{ij}v_{ij} - v_{ij}\nabla_{i}\nabla_{j} + 2 <\nabla f,\nabla h>
  + (R + \nabla f^2)({v \over 2} - h)]$$
  $$ = {d \over d f}F = {2 {\int (R + \nabla_{i}\nabla_{j} f)^2} \over 
  { - (R + \Delta f)}}\mathrm{dm} $$
  これらの方程式は8種類の微分幾何の次元多様体として、そして、これらの
  多様体は曲平面による双対性をも生成している。 
  そして、このガウスの曲平面は、大域的微分多様体と微分幾何の量子化から
  素因数を形成してもいる。
  These equation are eight differential geometry of dimension calyement.
  And these calyement equation excluded into pair of dimension surface.
  This surface of Gauss function are global differential manifold,
  and differential geometry of quantum level.
  $$ F \ge {d \over d f}\int \int{1 \over (x \log x)^2}dx_m 
  + {d \over d f}\int \int{1 \over (y \log y)^{1 \over 2}}dy_m $$

  微分幾何の量子化はオイラーの定数とガンマ関数が指数による連分数
  としての不変性として素因数を形成していて、このガウスの曲平面
  による量子力学における重力場理論は、ダランベルシアンの切断多様体
  がこの大域的切断多様体を付加してもいる。
  Differential geometry of quantum level constructed with Euler product
  and Gamma function being discatastrophed from continued fraction style.
  Gravity equation lend with varint of monotonicity of level expresented
  from gravity of letter varient formula.
  これらの方程式は基本群と大域的微分多様体をエスコートしていもいて、
  ヴェイユ予想がこのダランベルシアンの切断方程式たちから
  輸送のポートにもなっている。ベータ関数とガンマ関数がこれらのフォームラ
  の方程式を放出してもいて、結果、これらの方程式は広中平祐定理の複素
  多様体とグリーシャ教授によるペレルマン多様体からサポートされてもいる。
  この2名の教授は、一つは抽象理論をもう一方は具象理論を説明としている。
  These equation escourted into Global differential manifold and
  fundemental equation.
  Weil's theorem is imported from this equation in gravity of letters.
  Beta function and Gamma function are excluded with these formulas.
  These equation comontius from Heisuke Hironaga of complex manifold
  and Gresha professor of Perelman manifold.
  These two professos are one of abstract theorem and the other of
  visual manifold theorem.

\section{Hilbert manifold in Mebius space \\
        this element of Zeta function on integrate of fields}

 Hilbert manifold equals with Von Neumann manifold, and this fields is concluded with Glassman manifold,
the manifold is dualty of twister into created of surface built.
These fields is used for relativity theorem and quantum physics.
$$ ||ds^2|| = \lim_{x \to \infty} [\delta(x) \int \int \int \pi \left( \sum_{k=0}^{\infty}
		{{}^n\sqrt{p},x \over n} \right)
^{1 \over 2}d \tau]^{\mu\nu} $$
  This norm is component of fields in Hilbert manifold of space theorem rebuilt with
Mebius space in Gamma Function on Beta function fill into power of rout fundamental
group.
And this result with $AdS_5$ space time in Quantum caos of Minkowsky of manifold in
abel manifold.
Gravity of metric in non-commutative equation is Global differential equation
conbuilt with Kaluza-Klein space. Therefore this mechanism is $T^{\mu\nu}$ tensor
is equal with $R^{\mu\nu}$ tensor. And This moreover inspect with laplace operator
in stimulate with sign of differential operator.
Minus of zone in Add position of manifold is Volume of laplace equation rebuilt with
Gamma function equal with summative of manifold in Global differential equation,
this result with setminus of zone of add summative of manifold, and construct with
locality theorem straight with fundamental group in world line of surface, this
power is boson and fermison of cone in hyper function.

$$ V(\tau) = [f(x),g(x)] \times [f^{-1}(x), h(x)] $$
$$ \Gamma(p,q) = \int e^{-x} x^{1 - t}dx $$
$$ = \beta(p,q) $$
$$ = \pi (f(\chi,x),x) $$
$$ ||ds^2|| = \mathcal{O}(x)[(f(x) \circ g(x))^{\mu\nu}]dx^{\mu}dx^{\nu} $$
$$ = \lim_{x \to \infty} \sum_{k=0}^{\infty}a_k f^k $$
$$ G^{\mu\nu} = {\partial \over \partial f}\int [f(x)^{\mu\nu} \circ G(x)^{\mu\nu}dx^{\mu}dx^{\nu}]^{\mu\nu}dm $$
$$ = g_{\mu\nu}(x)dx^{\mu}dx^{\nu} - f(x)^{\mu\nu}dx^{\mu}dx^{\nu} $$
$$ [i \pi(\chi,x),f(x)] = i \pi f(x) - f(x) \pi(\chi,x) $$
$$ T^{\mu\nu} = (\lim_{x \to \infty} \sum_{k=0}^{\infty}\int \int [V(\tau) \circ S^{\mu\nu}(\chi,x)]dm)
^{\mu\nu}dx^{\mu}dx^{\nu} $$
$$ G^{\mu\nu} = R^{\mu\nu}T^{\mu\nu} $$
$$ \begin{vmatrix} D^m & dx \\
dx & \sigma^m \end{vmatrix} \begin{vmatrix} \cos \theta & - \sin \theta \\
\sin \theta & \cos \theta \end{vmatrix} \begin{vmatrix} x \\
y \end{vmatrix} = \begin{pmatrix} 1 & 0 \\ 
0 & - 1 \end{pmatrix}^{1 \over 2} $$
$$ \sigma^m \begin{bmatrix} \delta(x) & - 1 \\
1 & \epsilon(x) \end{bmatrix}^{1 \over 2} = \begin{pmatrix} i & 0 \\
0 & - i \end{pmatrix} $$
$$ V(M) = {\partial \over \partial f} ({}^N \int [f \smallsetminus M]^{\oplus N})^{\mu\nu} dx^{\mu}dx^{\nu} $$
$$ V(M) = \pi(2 \int \sin^2 dx)\oplus {d \over df}F^M dx_m $$
$$ \lim_{x \to \infty}\sum_{k=0}^{\infty} a_k f^k = \int (F(V)dx_m)^{\mu\nu}dx^{\mu}dx^{\nu} $$
$$ \bigoplus_{k=0}^{\infty}[ f \smallsetminus g] = \vee(M \wedge N) $$
$$ \pi_1 (M) = e^{-f 2 \int \sin^2 x dm} + O(N^{-1}) $$
$$ = [i \pi(\chi,x),f(x)] $$
$$ M \circ f(x) = e^{-f \int \sin x \cos x dx_m} + \log (O(N^{-1})) $$

 Non-Symmetry space time.
 $$ {d \over dL}V(\tau) = {d \over df}\int \int_M {1 \over (x \log x)^2}dx_m + {d \over df}\int \int_M
 {1 \over (y \log y)^{1 \over 2}}dy_m $$
 $$ \epsilon S(\nu) = \square_v \cdot {\partial \over \partial \chi}({}^5 \sqrt{\wedge g^2})d\chi $$
 Differential Volume in $AdS_5$ graviton of fundamental rout of group.
 $$ \wedge (F^m_t)^{''} = {1 \over 12}g_{ij}^2$$
 Quarks of other dimension.
 $$ \pi(V_{\tau}) = {e^{-(\sqrt{\pi \over 16}\log x)}}^{\delta} \times {1 \over (x \log x)} $$
 Universe of rout, Volume in expanding space time.
 $$ {d \over dt}(g_{ij})^2 = {1 \over 24}(F^m_t)^2 $$
 $$ m^2 = 2 \pi T \left({26 - D_n \over 24}\right) $$
 This quarks of mass in relativity theorem, and fourth of universe in three manifold of one dimension surface,
 and also this integrate of dimension in conbuilt of quarks.
 $$ g_{ij} \wedge \pi(\nu_{\tau}) = e^{-2 \pi T |\psi|}[\eta_{\mu\nu} + \bar{h}_{\mu\nu}(x)]dx^{\mu}dx^{\nu} 
 + T^2 d \psi^2 $$
 Out of rout in $AdS_5$ space time.
 $$ ||ds^2|| = g_{ij}\wedge \pi(\nu_v) $$
 $AdS_5$ norm is fourth of universe of power in three manifold out of rout.


  Sphere orbital cube is on the right of hartshorne conjecture by the right equation,
this integrate with fourth of piece on universe series, and this estimate with one of three manifold.
Also this is blackhole on category of symmetry of particles.
These equation conbuilt with planck volume and out of universe is phologram field, this field construct with
electric field and magnitic field stand with weak electric theorem.
This theorem is equals with abel manifold and $AdS_5$ space time.
Moreover this field is anti-brane and brane emerge with gravity and antigravity equation restructed with.
Zeta function is existed with Re pole of ${1 \over 2}$ constance reluctances.
This pole of constance is also existed with singularity of complex fields.
Von Neumann manifold of field is also this means.
$$ ||ds^2|| = \lim_{x \to \infty} [\delta(x)\int \int \int \pi \left( \sum_{k=0}^{\infty}{{}^n\sqrt{p},x \over n}\right)
^{1 \over 2}d \tau]^{\mu\nu} $$
$$ e^{i \theta} = \cos \theta + i \sin \theta $$
$$ e^{x \log x} = x^{{1 \over 2} + iy}, x \log x = \log (\cos \theta + i \sin \theta) $$
$$ = \log \cos \theta + i \log \sin \theta $$
This equation is developed with frobenius theorem activate with logment equation.
This theorem used to
$$ \log(\sin \theta + i \cos \theta)= \log(\sin \theta - i \cos \theta) $$
$$ \log \left({\sin \theta \over i \cos \theta}\right) = -2 R_{ij}, {d \over dt}g_{ij}(t) =  - 2 R_{ij} $$
$$ \mathcal{O}(x) = {\zeta(s) \over \sum\limits_{k=0}^{\infty}a_k f^k} $$
$$ \mathrm{Im} f = \mathrm{ker} f, \chi(x) = {\mathrm{ker}f \over \mathrm{Im}f} $$ $$ H(3) = 2, \nabla H(x) = 2, 
 \pi(x) = 0 $$
 This equation is world line, and this equation is non-integrate with relativity theorem of rout constance.

$$ [f(x)] = \infty, ||ds^2|| = \mathcal{O}(x)[\eta_{\mu\nu} + \bar{h}_{\mu\nu}(x)]dx^{\mu}dx^{\nu} + T^2 d^2 \psi $$
$$ T^2 d^2 \psi = [f(x)] , T^2 d^2 \psi = \lim_{x \to 1}\sum_{k=0}^{\infty}a_k f^k $$

 These equation is equals with M theorem. For this formula with zeta function into one of universe in four dimensions.
Sphere orbital cube is on the surface into $AdS_5$ space time, and this space time is Re pole and Im pole of 
constance reluctances.


  Gauss function is equals with Abel manifold and Seifert manifold. Moreover this function is infinite time,
and this energy is cover with finite of Abel manifold and Other dimension of seifert manifold on the surface.


  Dilaton and fifth dimension of seifert manifold emerge with quarks and Maxwell equation, this power
is from dimension to flow into energy of boson. 
This boson equals with Abel manifold and $AdS_5$ space time, and this space is created with Gauss function.


  Relativity theorem is composited with infinite on D-brane and finite on anti-D-brane, This space time
restructed with element of zeta function.
Between finite and infinite of dimension belong to space time system, estrald of space element have with 
infinite oneselves. Anti-D-brane have infinite themselves, and this comontend with fifth dimension of AdS5 
have for seifert manifold, this asperal manifold is non-move element of anti-D-brane and move element 
of D-brane satisfid with fifth dimension of seifert algebra liner.
Gauss function is own of this element in infinite element.
And also this function is abel manifold.
Infinite is coverd with finite dimension in fifth dimension of AdS5 space time.
Relativity theorem is this system of circustance nature equation.
AdS5 space time is out of time system and this system is belong to infinite mercy.
This hiercyent is endrol of memolite with geniue of element.
Genie have live of telomea endore in gravity accesorlity result.
AdS5 space time out of over this element begin with infinite assentance.
Every element acknowlege is imaginary equation before mass and spiritual envy.

  $$ T^2 d^2 \psi = \lim_{x \to 1}\sum_{k=0}^{\infty}a_k f^k $$
  $$ \lim_{x \to 1}\sum_{k=0}^{\infty}a_k f^k = [T^2 d^2 \psi] $$

  $$ {d \over dL}V(\tau)= {d \over df}\int \int_M({}^5\sqrt{x^2})d\Lambda 
  + {d \over df}\int \int_M {}^N({}^3\sqrt{x})^{\oplus N}d\Lambda $$
  $$ {}^M(\vee(\wedge f \circ g)^N)^{1 \over 2} = {d \over df}\int \int_M{1 \over (x \log x)^2}dx_m
  + {d \over df}\int \int_M{1 \over (y \log y)^{1 \over 2}}dy_m $$
  $$ ||ds^2|| = \mathcal{O}(x)[\eta_{\mu\nu} + \bar{h}_{\mu\nu}(x)]dx^{\mu}dx^{\nu}
  + T^2 d^2 \psi $$
  $$ \mathcal{O}(x) = e^{-2\pi T|\psi|} $$
  $$ G^{\mu\nu} = R_{\mu\nu} T^{\mu\nu} $$
  $$ = - {1 \over 2}\Lambda g_{ij}(x) + T^{\mu\nu} $$

  Fifth dimension of eight differential structure is integrate with one geometry element,
Four of universe also integrate to one universe, and this universe represented with symmetry formula.
Fifth universe also is represented with seifert manifold peices.
All after universe is constructed with six element of circuatance.
This aquire maniculate with quarks of being esperaled belong to.

\section{Imaginary equation in AdS5 space time create with \\
       dimension of symmetry}

  D-brane and anti-D-brane is composited with all of series universe emerged for one geometry of 
dimension, this gravity of power from D-brane and anti-D-brane emelite with ancestor.
Seifert manifold is on the ground of blackhole in whitehole of power pond of senseivility.
Six of element of quarks and universe of pieces is supersymmetry of mechanism resolved with

hyper symmetry of quarks constructed to emerge with darkmatter.
This darkmatter emerged with big-ban of heircyent in circumstance of phenomena.

  D-brane and anti-D-brane equations is
  $$ {d \over dL}V(\tau) = {d \over df}\int \int_M {1 \over (x \log x)^2}dx_m
  + {d \over df}\int \int_M {1 \over (y \log y)^{1 \over 2}}dy_m $$
  $$ C = \int \int {1 \over (x \log x)^2}dx_m $$

  Euler constance is quantum group theorem rebuild with projective space involved with.


   虚数方程式は、反重力に起因するフーリエ級数の励起を生成する。それは、人工知能を生み出す、
   ５次元時空にも、この虚数方程式は使われる。AdS5の次元空間は、反ド・シッター時空の
   D-braneとanti-D-braneのcomformal場を生み出す。ホログラフィー時空は、この量子起因によるものである。
   2次元曲面によるブラックホールは、ガンマ線バーストによる５次元時空の構造から観測される。
   空間の最小単位によるプランク定数は、宇宙の大域的微分多様体から導き出される、AdS5の次元空間の
   準同型写像を形成している。これは、最小単位から宇宙の大きさを導いている。最大最小の方程式は、
   相加相乗平均を形成している。時間と空間は、宇宙が生成したときから、宇宙の始まりと終わりを既に生み出している。
   宇宙と異次元から、ブラックホールとホワイトホールの力がわかり、反重力を見つけられる。オイラーの定数は、
   この量子定数からわかる。虚数の仕組みはこの量子スピンの産物である。
   オイラーの定数は、 この虚時空の斥力の現存である。それは、非対称性理論から導かれる。
   不確定性原理は、AdS5のブラックホールとホワイトホールを閉３次元多様体に統合する５次元時空から求められる。
   位置と運動エネルギーが、空間の最小単位であるプランク定数を宇宙全体にする微細構造定数からわかり、
   面積確定から、アーベル多様体を母関数に極限値として、ゼータ関数をこの母関数に不変式として、
   ２種類ずつにまとめる４種類の宇宙を形成する８種類のサーストンの幾何化予想から導き出される。
   この閉３次元多様体は、ミラー対称性を軸として、６種類の次元空間を一種類の異次元宇宙と同質ともしている。
   複素多様体による特異点解消理論は、この原理から求められる。この特異点解消理論は、２次元曲面を３次元多様体
   に展開していく、時空から生成される重力の密度を反重力と等しくしていく時間空間の４次元多様体と虚時空
   から求められる。
   ヒルベルト空間は、フォン・ノイマン多様体とグラスマン多様体をこのサーストンの幾何化予想を場の理論
   既定値として形成される。
   この空間は、ミンコフスキー時空とアーベル多様体全体を表している。そして、この空間は、球対称性を
   複素多様体を起点として、大域的トポロジーから、偏微分を作用素微分として時空間をカオスから
   ずらすと５次元多様体として成立している。
   これらより、３次元多様体に２次元射影空間が異次元空間として、AdS5空間を形成される。
   偏微分、全微分、線形微分、常微分、多重微分、部分積分、置換微分、大域的微分、単体分割、
   双対分割、同調、ホモロジー単体、コホモロジー単体、群論、基本群、複体、マイヤー・ビーートリス完全系列、
   ファン・カンペンの定理、層の理論、コホモルティズム単体、CW複体、ハウスドロフ空間、線形空間、位相空間、
   微分幾何構造、モースの定理、カタストロフの空間、ゼータ関数系列、球対称性理論、スピン幾何、
   ツイスター理論、双対被覆、多重連結空間、プランク定数、フォン・ノイマン多様体、グラスマン多様体、
   ヒルベルト空間、一般相対性理論、反ド・シッター時空、ラムダ項、D-brane,anti-D-brane,コンフォーマル場、
   ホログラム空間、ストリング理論、収率による商代数、ニュートンポテンシャルエネルギー、剛体力学、
   統計力学、熱力学、量子スピン、半導体、超伝導、ホイーストン・ブリッチ回路、非可換確率論、Connes理論、
   これら、演算子代数を形成している、微分・積分作用素が、ヒルベルト空間に存在している多様体の特質を
   全面に押し出して、いろいろな多様体と関数そして、群論を形作っている。


$$ \begin{vmatrix} D^m & dx \\
dx & \sigma^m \end{vmatrix} \begin{vmatrix} \cos \theta & - \sin \theta \\
\sin \theta & \cos \theta \end{vmatrix} \begin{vmatrix} x \\
y \end{vmatrix} = \begin{pmatrix} 1 & 0 \\ 
0 & - 1 \end{pmatrix}^{1 \over 2} $$
$$ \sigma^m \begin{bmatrix} \delta(x) & - 1 \\
1 & \epsilon(x) \end{bmatrix}^{1 \over 2} = \begin{pmatrix} i & 0 \\
0 & - i \end{pmatrix} $$
$$ (D^m,dx)\cdot (\cos \theta, \sin \theta) \times (dx, \partial^m)\cdot (\cos \theta, \sin \theta) $$

$$ \begin{pmatrix} 1 && 0 \\
   0 && i \end{pmatrix} \beta(x,\theta) \begin{pmatrix} x \\
   y \end{pmatrix} $$
   $$ = \begin{pmatrix} i && 0 \\
   0 && - 1 \end{pmatrix} $$
   $$ l = {\sqrt{\hbar G \over c^3}}, \sigma^m \cdot \begin{pmatrix} \delta(x) && - 1 \\
   1 && \epsilon (x) \end{pmatrix}^{1 \over 2} = \begin{pmatrix} i && 0 \\
   0 && - i \end{pmatrix} $$
   String theorem also imaginary equation of pole with the fields of manifold.
   $$ \left({\partial \over \partial \tau}f(x,y,z)\right)^{3'} = A^{\mu\nu} $$
   $$ {d \over dt}g_{ij}(t) = - 2 R_{ij} $$
   Rich equation is all of top formula in integrate theorem.
      $$ = (D^m,dx) \cdot (\cos \theta, \sin \theta) \times (dx, \partial^m)\cdot (\cos \theta, \sin \theta) $$
      Dalanverle equation equals with abel manifold.






\title{慣性の法則と、回転体による反重力 \\
タイミングとともに、別ベクトルとして重力を加える法則}
\author{Masaaki Yamaguchi}


   慣性の法則の慣性力が働くときと同時に、
   前もって、タイミングを計り、慣性力とともに、
   逆向きのベクトルとして、慣性力の力と同じエネルギーの
   重力を生成して、慣性力の影響を無くす。この力を反重力という。
   慣性力と逆向きのベクトルとして、反重力として、
   慣性力に加えて、別のベクトルとする。

   反重力の本質は、重力のベクトルの違いであり、
   核エネルギーが特殊相対性理論と一般相対性理論を使って、
   放射性物質吸収体に、この原理を使い、核エネルギーをリサイクルしている。
   これを宇宙人は知っていて、UFOにこの原理を使っている。

   Higgs場は、重力であり、周りを反重力が覆っている。
   一般相対性理論のベクトルの向きを変えた力の装置が、
   ラムダ・ドライバであり、反重力である。

   重力は熱エネルギーである。ベクトルの向きが違う。
   反重力は吸収体である。
   慣性の法則が働いているときに、これが手に入ったが、離れたになり、
   重力をベクトルの向きを変えて、同じエネルギーを加えた場合、
   量子コンピュータが、そのときの差分を使って、反重力が慣性の法則
   を変えて、宙に浮くUFOにしている。 

   まとめると、慣性の法則のときに、重力のベクトルを違う向きとして、
   加えると、手に入ったが、離れたになり、これが反重力である。 

   UFOが現れる前に、ヘリコプターが現れて、そのあとに、UFOが実際に現れる。
   慣性の法則のしばらくして、反重力として、UFOとして、
   宇宙人が人に教えている。

   反重力の生成は、回転体の電磁気力からくるちからであり、ローレンツ力と
   同じでもあり、機体の内部で生成される力であり、そのために、放射性
   物質吸収体を結晶石として、このエネルギーからくる放射力を吸収する
   ために、結晶石に反重力の放射線を吸収する。その上に、外部の
   慣性力を内部の反重力で、手に入ったら、離れたをしている。


\title{反重力発生装置の各構成物質}
\author{Masaaki Yamaguchi}


  原子振動子　セシウム Cs


  形状記憶合金 (Fe・Co60・Pt・Al)$H_2 SO_4, Al_2 O_3$


  反重力発生器 Pr(パラジウム合金)電磁場生成He,$H_2$Mg,Al(摩擦熱)


  結晶石 Co60,Pt,Prの反陽子


  緩衝剤 慣性力感知器 有機化合物とS化合物(シクロアルカン$C_n H_{2n}$
  方位と位置探知機にも使う


  量子コンピュータの量子素子 Pt,Ag,Au（電子と電流のエネルギー経路)半導体に
  使う
  回路の合成演算子 


  $$ x^y = {1 \over {y^x}} \to $$

  $$ \pi(\chi,x) = [i\pi(\chi,x),f(x)] $$
  $$ {d \over df}F = F^{{f}^{'}} $$


  ディスプレイの電磁迷彩 (Al,Mg) $\to$ S合成子 $Si O_2 $
  プリズム C
\title{特殊相対性理論の虚数回転による多様体積分と、\\
それによる一般相対性理論の再構築理論}
\author{Masaaki Yamaguchi}
   
  \newcommand{\oiint}{%
	   \begin{picture}(15,15)
		   \put(0,0){
			   $ \scalebox{1}[2]{$\iint$} $}
		   \put(2,2){$\bigcirc$}
	   \end{picture} %
	   }


 $$ {\partial \over \partial f}F(x) = \int \int \text{cohom}D_k(x)[I_m]  $$
 Cohomology element is constructed with isotopy of D-brane,
 and this isotopy of symmetry structure is sheaf of manifold,
 more over this brane of element is double integrate of projection.
 $$ {\partial \over \partial f}F = {{}^t \scalebox{1}[2]{\variint}}
 \text{cohom}D_k(x)^{\ll p} $$

  $$  {\partial \over \partial f}F = {{}^t \scalebox{1}[2]{\variint}}
 \text{cohom}D_k(x)^{\ll p} = 
\scalebox{1}[2]{\alearletter}$$ 
   $$ \nabla_{i} \nabla_{j} \int f(x) d \eta = {\partial^2 \over {\partial x
   \partial y}}\int \scalebox{1}[2]{\alearletter} d \eta $$
   一般相対性理論の加群分解が偏微分方程式と同じく、
   特殊相対性理論の多様体積分の虚数回転体がベータ関数となる。
   ほとんどの回転体の体積が、係数と冪乗での回転体として、ベータ関数
   と言える。
   $$ = R^{{\mu\nu}^{'}} + {1 \over 2}\Lambda g_{ij}^{'} = \int 
   \left(i {{v \over {
	   \sqrt{1 - 
   {{{({v \over t})}}^{2}}}}}} + {v \over {\sqrt{1 - 
   {{{({v \over t})}}^{2}}}}}\right)\mathrm{dvol} $$
   $$ = {d \over df}\int \int{1 \over {(x \log x)^2}}dx_m +
   {d \over df}\int \int {1 \over {(y \log y)^{1 \over 2}}}dy_m $$ 
    この大域的積分多様体が大域的微分多様体の反重力と重力方程式
    で表せられて、
   $$ = \int C dx_m = \int \kappa T^{\mu\nu} \mathrm dx_m = T^{{\mu\nu}{^{T^
   {\mu\nu}}}^{'}}$$
   オイラーの定数の大域的積分多様体が、一般相対性理論の大域的積分多様体
   であり、エントロピー不変式で表せられる。

  $$ \scalebox{1}[2]{\alearletter} = 
  {{{{8 \pi G} \over {c^4}}T^{\mu\nu}}/ \log x} $$
  $$ {}^{t}\scalebox{1}[2]{\variiint}\mathrm{cohom D_{\chi}[I_m]}$$
  $$ = \oint {(px^n + qx + r)^{\nabla l}} $$
  $$ {d \over d l}L(x,y) = 2 \int ||\sin 2x||^{2}d \tau $$
  $$ {d \over d \gamma}\Gamma $$
 $$ ||ds^2|| = \lim_{x \to \infty} [\delta(x) \int \int \int \pi \left( \sum_{k=0}^{\infty}
		{{}^n\sqrt{p},x \over n} \right)
^{1 \over 2}d \tau]^{\mu\nu} $$

 $$\oiint  \cong  ||ds^2|| = \lim_{x \to \infty} 
 [\delta(x) \int \int \int \pi \left( \sum_{k=0}^{\infty}
		{{}^n\sqrt{p},x \over n} 
		\right)^{1 \over 2}d \tau]^{\mu\nu} $$
 
 $$ e^{-2\pi T||\psi||}[\eta + \bar{h}]dx^{\mu}dx^{\nu} + T^2 d^2 \psi $$
 $$ {{-16 \pi G} \over c^4}T^{\mu\nu}/ \log x =  
 \scalebox{1}[2]{\alearletter} $$
 $${{-16 \pi G} \over c^4}T^{\mu\nu}/ e^{-2\pi T||\psi||} $$
 $$ = 4 \pi G \rho $$
 $$ {\partial \over {\partial x \partial y}} = \nabla_{i} \nabla_{j} $$ 
 $$ \square \iiint  = {}^{t}\scalebox{1}[2]{\variiint}\ $$
 $$ {\partial \over {\partial x}} \iiint = \nabla_{i}\nabla_{j} \int
 \nabla f(x) d \eta $$
 $$ {\partial \over {\partial f}}\iiint = \square \iiint $$

 $$ \reflectbox{$\Gamma$} | \Gamma, \reflectbox{B} | B $$
 $$ \reflectbox{E} | E, \reflectbox{C} | C $$
 $$ \reflectbox{F} | F, \reflectbox{$\beta$} | \beta $$
 $$ \reflectbox{D} | D, {}^{t}\scalebox{1}[2]{\variint} \cong 
 \bigoplus D^{\bigoplus L} $$
 $$ \square +  \scalebox{1}[2]{\alearletter} = \emptyset $$
 $$  \scalebox{1}[2]{\alearletter} | \emptyset = \square $$
 $$ {\begin{pmatrix} 1 & 0 \\
 0 & -i \end{pmatrix}}^{1 \over 2} \Biggl| {\begin{pmatrix} 1 & 0 \\
 0 & -i \end{pmatrix}}$$
 $$ x^{1 \over 2} \cong x, \emptyset^{1 \over 2} = \emptyset$$
 
    自分がいると、方位が決まる。そうしたら、空間ができる。北と南が出来る。
    これは、磁気単極子が相対空間での根本的なアイデアらしい。
    無になると関係性がなくなり、磁気単極子もある上に、全て無でもある。
    無になると十二支縁起になる。宇宙と異次元の外が無でもある。
    運動量が決まると位置が決まる。位置だけだと宇宙が不確定性原理
    であり、異次元が合わさると運動量も位置も決まる。
    無になると位置も運動量も同時にわかる。
    このアイデアが3次元多様体のエントロピー値である。
    
    $$ \int^{\ll n} ||ds^2|| = \int\int\int 8 \pi G \left({p \over c^3} + 
    {V \over S}\right)dxdydz $$
    $$ = {\partial \over \partial f_n}{||ds^2||} $$


 $$ {\partial \over \partial f}F(x) = \int \int \text{cohom}D_k(x)[I_m]  $$
 Cohomology element is constructed with isotopy of D-brane,
 and this isotopy of symmetry structure is sheaf of manifold,
 more over this brane of element is double integrate of projection.
 $$ {\partial \over \partial f}F = {{}^t \scalebox{1}[2]{\variint}}
 \text{cohom}D_k(x)^{\ll p} $$
 And global partial defferential equation is integrate of cut in cohomolgy,
 and this step of defferential D-brane of sheap value is varint of equals,
 inverse of horizen of cute of section value is reverse of integrate of 
 operators. Three dimension of varint of manifolds in operator of sheaps,
 and this step of times of value is equal with D-brane of equations.
 Global differential equation and integrate equation is operator of 
 global scales of horizen cut and add of cut equation are routs.

 $$ \oint r dx_m = \mathcal{O}(x,y) $$

 $$ {{}^t \scalebox{1}[2]{\variiint}}_{{D(\chi,x)}} 
 \text{Hom}[D^2 \psi]^{\ll p} \cong 
 \text{vol}\left({V \over S}\right) $$

 $$ {\partial \over \partial f}F(x) = {{}^t \scalebox{1}[2]{\variint}} 
 \text{cohom}D_k(x)^{\ll p} $$

 $$ {d \over df}F = (F)^{f^{'}}, \int F dx_m = (F)^{f}  $$
 $$ {d \over df}\int \int{1 \over (x \log x)^2}dx_m =
 \left(\int \int{1 \over (x \log x)^2}dx_m \right)^{f^{'}}  $$
 $$ f^{'} = (2x (\log x + \log (x + 1))) $$
  $$  \int \int F dx_m = \left({1 \over (x \log x)^2}\right)^{
	  (\int 2x (\log x + \log (x + 1))dx)} $$ 
   $$= e^{-f} $$
  $$ {d \over df}F = \left({1 \over (x \log x)^2}\right)^{
	  (2x (\log x + \log (x + 1))) } $$  
  $$= e^{f} $$
  $$ \log (x \log x) \ge 2(y \log y)^{1 \over 2} $$

  $$ \log (x \log x) \ge 2(\sqrt{y \log y}) $$

 
    $$ ||ds^2|| = 8 \pi G \left({p \over c^3} + {V \over S}\right) $$
    $$ \int ||ds^2||dx_m  = \int 8 \pi G \left({p \over c^3} 
    + {V \over S}\right)\mathrm{dvol}  $$
    $$ \int ||ds^2||dx_m  = \int 8 \pi G \left({p \over c^3} 
    + {V \over S}\right)dy_m $$
     $$ \int ||ds^2||dx_m  = \int {1 \over {8 \pi G \left({p \over c^3} 
    + {V \over S}\right)}}dx_m $$
    $$ {d \over df}F = m(x), \bigoplus \left({i \hbar}^{\nabla}\right)^{
	    \oplus L} $$
    $$ = \nabla_{i} \nabla_{j} \int \nabla f(x)d \eta $$

    $$ dx_m = {y \over {\log x}}, dy_m = {x \over \log x} $$
    $$ e^{-f}dV = dy_m = \mathrm{dvol} $$
    偏微分は加群分解と同じ計算式に行き着く。

   宇宙と異次元の誤差関数のエネルギー、
    $AdS_5$多様体がベータ関数となる値の列が、異次元への扉となっている。
    $$ \beta(p,q) = \text{誤差関数} + \text{Abel多様体} $$
    $$ = \text{$AdS_5$多様体} $$
    $$ = {d \over df}F + \int C dx_m = \int \Gamma(\gamma)^{'}dx_m $$
    ここで、アーベル多様体はEuler productである。
    ベータ関数の数列がわかると、ゼータ関数は無であるというのが、
    どういうことかが、物、物体に影ができて、ものが瞑想と同じであり、
    これから、風景がベータ関数の数値列に見えるらしい。
     この大域的微分多様体のガンマ関数が、複素力学系のマンデルブロ集合の
    プリズムと同じ構造の見方らしい。

    $$ ||ds^2|| = e^{-2\pi G ||\psi||}[\eta + \bar{h}(x)]dx^{\mu}dx^{\nu}
    + T^2 d^2 \psi $$
    $$ \beta(p,q) = \text{誤差関数} + \text{Abel多様体} $$
    $$ \int \mathrm{dvol} = \square \psi $$
    $$ \int \nabla \psi^2 d \nabla \psi = \square \psi $$
    $$ \text{expanding of universe} = \text{exist of value} $$
    $$ = \log (x \log x) = \square \psi $$
    $$ \text{freeze out of universe} = \text{reality of value} $$
    $$ = (y \log y)^{1 \over 2} = \nabla \psi $$

    All of value is constance of entropy, universe is freeze out
    constant, and other dimension is expanding into fifth dimension of
    inner.


      $$ x^n + y^n = z^n, \beta(p,q) = x^n + y^n - \delta(x) = z^n - \delta
      (x) $$
      $$ {d \over dt}g_{ij} = - 2 R_{ij}, {\cos x \over (\cos x)^{'}}
      \cdot (\sin x)^{'} = z_n, z^n = - 2 e^{x \log x} $$

      $$ z = e^{-f} + e^{f} - y $$
      $$ \beta(p,q) = e^{-f} + e^{f} $$
      相対性は、暗号解読と同じ仕組みの数式を表している。
      ここで言うと、yが暗号値である。テェックディジットと同じ仕組みを有
      している。

      $$ \square x = \int {f(x) \over {\nabla(R^{+}\cap E^{+})}}d \square x$$
      $$= \int {{\Delta f(x) \circ E^{+}} \over {\nabla(R^{+}\cap E^{+})}}
      \square x $$
      $$ \square x = \int {d(R \nabla E^{+}) \over {\nabla(R^{+} \cap E^{+})}}
      d \square x $$
      $$ = \int {{\nabla_{i} \nabla_{j} (R + E^{+}} \over {\nabla(R^{+}
      \cap  E^{+})}}d \square x $$
      $$ x^n + y^n = z^n $$
      $$ \mathrm{exp}(\nabla(R^{+}\cap E^{+}),\Delta(C \supset R)) $$
      $$ = \pi (R_1 \subset \nabla E^{+}) = \mathrm{rot}(E_1, 
      \mathrm{div}E_2) $$
      $$ xf(x) = F(x), s\Gamma(s) = \Gamma(s+1) $$
      $$ Q \nabla C^{+} = {d \over df}F(x) \nabla \int \delta(s)f(x)dx $$
      $$ E^{+}\nabla f = {{e^{x \log x} \nabla n! f(x)} \over {E(x)}} $$
      $$ {d \over df}F ~ F^{{f}^{'}} = e^{x \log x} $$
      $$ (C^{\nabla})^{\oplus Q} = e^{x \log x} $$
      $$ {d \over dt}g_{ij} = - 2 R_{ij}, {\cos x \over (\cos x)^{'}}
      \cdot (\sin x)^{'} = z_n, z^n = - 2 e^{x \log x} $$

      $$ R \nabla E^{+} = f(x) \nabla e^{x \log x}, {d \over df}F = F^{
	      {f}^{'}} = e^{x \log x}  $$
      南半球と単体（実数）の共通集合の偏微分した変数をどのようなF(x)かを
      $$ \int \delta(x)f(x)dx $$
      と同じく、単体積分した積分、共通集合の偏微分をどのくらいの微分変数
      を
      $$ \int \nabla \psi^2 d \nabla \psi = \square \psi $$
      と同じ、
      $$ \int dx = x + \text{C(Cは積分定数) }$$
      と原理は同じである。


      Beta function is,
      $$ \beta(p,q) = \int x^{1 - t} (1 - x)^{t}dx = \int t^{x}(1 - t)^{x - 1}
      dt $$
      $$ 0 \le y \le 1, \int^{1}_{0} x^{10}(1 - x)^{20}dx = B(11,21) $$
      $$ = {{\Gamma(11)\Gamma(21)} \over \Gamma(32)} = {{10! 20!} \over {
	      31!}} = {1 \over {931395465}} $$
       $$ {1 \over {931395465}} = {1 \over 9} = {1 \over {1 - x}} $$
	 $$     = {1 \over {1 - z}} =
	      \sum_{k=0}^{\infty}z^k = {1 \over {1 + z^2}} = \sum_{
		      k=0}^{\infty} (-1)^k z^{2k} $$
      $$ f(x) = \sum_{k=0}^{\infty} a_k z^k $$
      $$ {{d^n y} \over dx^n} = n! y^{n+1} $$
      $$ f^{(0)}(0) = n!f(0)^{n+1} = n! $$
      $$ f(x) \cong \sum_{k=0}^{\infty}x^n = {1 \over {1 - x}} $$
      In example script is,
      $$ {{dy} \over dx} = y^2, {1 \over y^2}\cdot {dy \over dx} = 1 $$
      $$ \int {1 \over y^2}{dy \over dx}dx = \int{dy \over y^2} = - {1 \over y
      } $$
      $$ - {1 \over y} = x - C, y = {1 \over {C - x}} $$
      $$ \exists x = 0, y = 1 $$
      $$\text{in first value condition compute with} $$
      $$ \text{result, C consumer sartified,} $$
      $$ y = {1 \over {1 - x}} $$
      This value result is concluded with native function from Abel manifold.

      Abelian enterstane with native funtion meltrage from Daranvelcian
      equation on Beta function.
 
       Rotate with pole of dimension is converted from 
       other sequence of element,this atmosphere of vacume 
       from being emerged with quarks of being constructed
with Higgs field, and this created from energy of zero dimension are begun with
universe and other dimension started from darkmatter that represented with 
big-ban.
Therefore, this element of circumstance have with quarks of twelve lake with
atoms.

  $$ (dx, \partial x)\cdot (\epsilon x, \delta x) = \begin{pmatrix}
	  1 & 0 \\
  0 & - 1 \end{pmatrix}^{1 \over 2}  $$
  $$ {d \over df} F = m(x) $$

 And this phenomounen is super symmetry theorem built with quarks of element, 
also this chemistry of mechanism estimate from physics of operatsion.
These response of mechanism chain of geometry into space being emerged with
creature and univse of existing of combination.

  This converted with dimension of element also emerged from imaginary and 
reality of pole's space.

  And this pole of transport of dimension belong with vector of time has 
with one rout of sequecence. Quanum physics also belong with other vector 
of time has with imaginary rout of sequecense.

 This sequence of being estimated with non fluer of time, and this space 
of element have with gravity and antigravity of power.
Other vecor of time is antigravity rout of sequecense.

 Weak electric theory is estimate from time has with one rout of sequecense,
this topology of chain is resulted from time of rout ways.

      $$ \square ({{\sigma_1 + \sigma_2} \over 2 }) = [3 \pi (\chi,x) \circ f(x),
   \sigma(x)] $$
   $$ \square \psi = 8 \pi G T^{\mu\nu} $$
   $$ {d \over dt}g_{ij}(t) = - 2 R_{ij} $$
   $$ \nabla \psi^2 = 4 \pi G \rho $$
   $$ \square (\sigma_1 + \sigma_2) = [6\pi(\chi,x)\circ f(x), \sigma(x)] $$
   $$ = [i \pi(\chi,x),f(x)] $$
   $$ \square \psi = [12 \pi (\chi,x)\circ f(x), \sigma(x)] $$
   $$ {d \over df} F = \int e^{-f}[- \Delta v + R_{ij} v_{ij}+ 
   \nabla_i \nabla_j f  + v \nabla_i \nabla_j + 2<f,h> + (R + \nabla f)
   ({v \over 2} - h) ] $$
   $$ =[i \pi(\chi,x), f(x)] $$

   These equation is reminded time pass rout of one rout way of forms,
   and this rout of time ways which go for system from future and past.
   Therefore, this resulted system of time mechanism is one true flow
   that weak electric theorem oneselves.
   and moreover, one rout time way of forms is reverse with antigravity of
   time system. This also spectrum focus is true that Maxwell theorem and
   strong boson unite with antigravity, this unite is essense on the contrary
   from weak electric theorem, this theorem called for time rout forms is
   strong electric theorem.
   This two theorem is united with quantum physics that no time flow system.

   $$ f^{-1}(x)xf(x) = 1,H_m = E_m \times K_m $$


   The non-commutative theorem is constructed from world line surface that
   this complex manifold estimate with rolanz atractor, and this string
   theorem have with one world of universe mate six quarks and other world of
   dimension mate other element of six quarks.
   These particle quarks is built with super symmetry space of dimension.

   $$ i = (1,0)\cdot (1,0), e^{i\theta} = \cos \theta + i \sin \theta $$
   $$ e^{-i \theta} = \cos \theta - i\sin \theta $$
   $$ \sin \theta = {e^{i \theta} - e^{-i \theta} \over 2i} $$
   $$ \sin i \theta = {e^{-\theta} - e^{\theta} \over 2} $$
   $$ \pi (\chi,x) = \cos \theta + i \sin \theta $$

   In this equations, two dimension redestructed into three dimension,
   this destroy of reconstructed way is append with fifth dimension.
   This deconstructed way of redestructe is arround of universe attached
   with three dimension, this over cover call into fifth dimension.

   $$ R(-\alpha) = \begin{pmatrix} \cos \alpha & \sin \alpha \\
   - \sin \alpha & \cos \alpha \end{pmatrix} $$
   $$ R(\alpha) = \begin{pmatrix} \cos \alpha & - \sin \alpha \\
   \sin \alpha & \cos \alpha \end{pmatrix} $$
   $$ R(\alpha)M R(-\alpha) = 
   \begin{pmatrix} \cos \alpha & - \sin \alpha \\
   \sin \alpha & \cos \alpha \end{pmatrix} \begin{pmatrix} 1 & 0 \\
   0 & - 1 \end{pmatrix} \begin{pmatrix} \cos \alpha & \sin \alpha \\
   - \sin \alpha & \cos \alpha \end{pmatrix} $$
   $$ = \begin{pmatrix} \cos^2 \alpha - \sin^2 \alpha & 2 \sin \alpha \cos 
	   \alpha \\
	   2 \sin \alpha \cos \alpha & - \cos^2 \alpha + \sin^2 \alpha
   \end{pmatrix} $$
   $$ = \begin{pmatrix} \cos 2 \alpha & \sin 2 \alpha \\
   \sin 2 \alpha & - \cos 2 \alpha \end{pmatrix} $$

   $$ \begin{pmatrix} \cos \alpha & - 1 \\
	   1 & - \sin \alpha \end{pmatrix} \begin{pmatrix} \cos \alpha \\ 
		   \sin \alpha
   \end{pmatrix} = \begin{pmatrix} 1 & 0 \\
   0 & - 1 \end{pmatrix} $$
   $$ \begin{pmatrix} \cos \alpha & \sin \alpha \\
   \sin \alpha & - \cos \alpha \end{pmatrix} \begin{pmatrix}
   x \\ y \end{pmatrix} = \begin{pmatrix} 1 & 0 \\
   0 & - 1 \end{pmatrix} $$
   $$ \lim_{\theta \to 0} \begin{pmatrix} \sin \theta \\ \cos \theta 
   \end{pmatrix} \begin{pmatrix} \theta & 1 \\
   1 & \theta \end{pmatrix} \begin{pmatrix} \cos \theta \\ \sin \theta
	   \end{pmatrix} = \begin{pmatrix} 1 & 0 \\ 0 & - 1 \end{pmatrix} $$
		   $$ f^{-1}xf(x) = 1 $$
		   $$ (\log x)^{'} = {1 \over x}, x^n + y^n = z^n  $$
		   $$ x^n = - y^n + c, nx^{n-1} = -n y^{n-1}y^{'} $$
		   $$ y^{'} = {{n x^{n-1}} \over {n y^{n-1}}} $$
		   $$ = {x^{n-1} \over y^{n-1}}= - {y \over x}\cdot
		   ({x \over y})^n $$
		   $$ - {\cos x \over (\cos x)^{'}}(\sin x)^{'} = z_n $$
		   $$ z^n = - 2 e^{x \log x} $$
		   $$ \lim_{x \to \infty} f(x) = a, \lim_{y \to \infty}
		   f(y) = b, \lim_{x,y \to \infty}\{f(x) + f(y)\} = a + b  $$
		   $$ \lim_{z \to 0}f(z) = c, \delta \int z^n = {d \over dV}
		   x^3 $$
		   $$ \lim_{x \to \infty} = c - \lim_{y \to \infty}f(y)
		   \lim_{x \to \infty} f(x) + \lim_{y \to \infty}f(y) = 
		   \lim_{z \to \infty}f(z) $$
		   $$ z^n = \cos n \theta + i \sin n \theta $$
		   $$ = - 2 e^{x \log x} $$

  These equation is gravity and antigravity equation. 

	   $$ {d \over d\sigma}\begin{bmatrix}{(\sigma_1 + \sigma_2) 
	   \over 2}\end{bmatrix}  $$
	   $$ = \sigma(\upharpoonleft \downharpoonleft) + \sigma(
	   \upuparrows)+ \sigma(\downdownarrows)+ 
	   \sigma(\leftleftarrows)
	   + \sigma(\rightrightarrows) $$
	   $$ {d \over df}{\int \int {1 \over (x \log x)^2}}dx_m
	   = \sigma(\upharpoonright) $$
	   $$ {d \over df}{\int \int {1 \over (y \log y)^{1 \over2}}}dy_m
	   = \sigma(\downharpoonright) $$
	   $$ \sigma(\leftarrowtail) + \sigma(\rightarrowtail)
	   = \int e^{-f}[-\Delta v + R_{ij}v_{ij} + \nabla_{i}\nabla_{j}v
	   + v\nabla_{i}\nabla_{j} + 2 <f,h> + (R + \nabla f)(v - 
	   {h \over 2})]
	   $$
	   $$ \sigma(\downharpoonleft) = \sigma(\upharpoonright 
	   \downharpoonright + \upuparrows 
	   + \rightrightarrows)$$
	   $$ \sigma(\upharpoonleft) =  \sigma(\upharpoonright 
	   \downharpoonright + \downdownarrows 
	   + \leftleftarrows) $$
         
	   $$ \text{weak electric theorem} = \sigma(\upharpoonleft) $$
	   $$ \text{strong electric theorem} = \sigma(\downharpoonleft) $$

	   These equation is represented with topology of string model,
	   and weak electric theorem is constructed with Maxwell theorem 
	   and weak boson, gravity that estimate with this three power united.
	   Moreover, strong boson and Maxwell theorem, antigravity that
	   also estimate with this three power united.
	   This two united power is integrated from gravity and antigravity.
	   Then this united power is zeta function.

  Euler equation represet form,
  $$ C = \int {1 \over x^s}dx - \log x $$
  $$ = \int \left(\int {1 \over x^s}dx - \log x \right)\mathrm{dvol} $$
  $$ = \int \int{1 \over x^s}\mathrm{dvol} - \int  \log x \mathrm{dvol} $$
  This equation componet of $ \text{Higgs equation + Euler equation =
  Zeta function} $. This represent expression is Rich flow equation.
  This equation assent from Differential geometry in quantum level,
  Gamma and Beta function, Gauss surface, Higgs field, Einsetin tensor,
  Euler constance and so on.
    This equation is add and average of possibility, and universe, atom
  more over a certain theory, this universe oneselves is environment
  seek, other distance is possibility equation, however, other dimension
  cover demand with not only norm of non liner in three manifold
  but also non possibility in complete of environment equation.
    $$ = {d \over df} \int \int{1 \over (y \log y)^{1 \over 2}}dy_m 
  - \int e^{-x}x^{1 - t} \log x dx_m $$
  $$ \int x^{1 - t}e^{-x}dV = \int x^{1 - t}dm, \int x^{1 - t}e^{-x}dV
  = \int x^{1 - t}\mathrm{dvol}, f = \gamma = \Gamma^{'} = \int 
  e^{-x}x^{1 - t} \log x dx $$
  $$ = {d \over d \gamma}\Gamma^{-1} - (\gamma)^{\gamma^{'}} $$
  $$ = e^{-f} - e^{f} $$
  $$ = {d \over df}\int \int{1 \over (y \log y)^{1 \over 2}}dy_m
  - {d \over df}\int\int {1 \over (x \log x)^2}dx_m $$
  $$ = 2 \cos (i x \log x) $$
  $$  {d \over df}\int\int{1 \over (y \log y)^{1 \over 2}}dy_m
  + {d \over df}\int\int{1 \over (x \log x)^2}dx_m $$
  $$ = {d \over df}F = 2 i \sin (i x \log x) $$
  $$ {d \over df}F + \int C dx_m  
  = 2 (\cos( i x \log x) + i \sin(i x \log x)) $$
  $$ = 2 e^{- \theta} = 2 e^{- i x \log x} $$
  $$ = {d \over dt}g_{ij}(t) = - 2 R_{ij} $$
  $$ \log (x \log x) \ge 2 \sqrt{y \log y} $$
  $$ \log(x \log x) = \log x + \log \log x $$
  This universe and other dimension on blackhole and whitehole
  compont with universe and atom in movement and point of 
  environment value by a certain theory of complete environment 
  equation. And this three manifold of possibility equation 
  are able to research in both value formula.
  This envirnment of researchable reason is that surround of universe
  is setting to value by already in consented to being researched,
  then both value researchabled.
  And this three manifold equation is insert with Kaluza-Klein theory
  and Heisenbelg equation.
    $$ \nabla \psi^2 = 8 \pi G\left({p \over c^3} + {V \over S}\right) $$
  $$ \nabla \psi^2 = 8 \pi G \hbar + 8 \pi {V \over S} $$
  This energe of entropy value consist with whitehole of atom
  and universe in blackhole value, and this universe in balckhole 
  represent with Higgs filed, atom in whitehole consist with plack scale.
    $$ \square_v = 2 \sqrt{2 \pi G \hbar} + 2 \sqrt{2 \pi {V \over S}} $$
  This entropy value also consist with atom in balckole.
    $$ \sqrt{2 \pi T} = 2 \sqrt{2 \pi{V \over S}} $$

  These equation more also consist with universe being non gravity formula.
    $$ 2 \cos (i x \log x) + 2 i \sin(i x \log x) = 2 e^{-f} $$
  $$ 2 \sqrt{2 \pi {V \over S}} = {d \over df} 
  \int \int{1 \over (x \log x)}dx_m$$
  And also these equation consist with universe and other dimension in
  warped passeage in rout of integral.

   $$ e^{-2\pi ||psi|}[\eta_{\mu\nu} + \hbar{x}]dx^{\mu}dx^{\nu} 
  +  T^2 d^2 \psi $$
  $$ = e^{-f} + e^f = (u + d) + c, (w + s) + b = -e^{-f} + e^f $$
  この式は、時空にどれだけの原子が凝縮しているかを表していて、
  その上に、この逆数は、密度エネルギーをも示している。
  空間の体積にに原子を商代数にすると、濃度にもなる。
   逆数は、$\text{a=原子のエネルギー}$が全空間を1とすると
  $\text{$1 \over {y \over 
  x}$}$すると、$x \over y $は密度になる。　
  $\text{{x=原子},y=全空間、個数は}$$\text{$y \over x$}$ 
  $$ \square = \text{原子} \to {{[\eta_{\mu\nu} + \hbar(x)]dx^{\mu}dx^{\nu}
  } \over {e^{2 \pi T ||\psi||} }}$$
  $$ ||ds^2|| = e^{2 \pi T||\psi||}\left([\eta_{\mu\nu} + \hbar(x)]d^{\mu}dx^
  {\nu}\right)^{-1} + \left(T^2 d^2 \psi \right)^{-1} $$
  $$||ds^2|| = e^{-2 \pi T||\psi||}[\eta_{\mu\nu} + \hbar(x)]d^{\mu}dx^
  {\nu}] + T^2 d^2 \psi  $$
  $$ (u + d) + c = e^{-f} + e^f, (w + s) + b = e^{-f} + e^{f} $$
  $$ R_{ij}^{'} = - R_{ij} $$

  この$ AdS_5 $多様体の $e^{-2 \pi T ||\psi||} $は原子間距離を表していて、
  $ [\eta + \bar{h}(x)]d^{\mu}dx^{\nu} $で宇宙のD-braneの構造を表している。
  この数式が$T^2 d^2 \psi $とアーベル多様体が包み込んでいる。
  原子が回転するのと、ニュートンリングが回転する宇宙は同じ速さで回転する。　

  $$||ds^2|| = \square + \rho , ||ds^2|| = \nabla \Psi^2 + \psi $$
  $$ \eta_{\mu\nu} + \bar{h}(x) $$
  $$ \text{ploximetry} = \Delta x \Delta p - \Delta p \Delta x + 
  \delta(p,x)  $$
  $$ ||ds^2|| = e^{-2 \pi T ||\psi||} + [\eta_{\mu\nu}  + \bar{h}(x)]
  dx^{\mu\nu} + T^2 d^2 \psi $$
  $$ {d \over df}F = e^f + e^{-f}, {d \over df}\int C dx_m = e^f - e^{-f} $$
  $$ R_{ij}^{'} = - R_{ij} $$
  宇宙と異次元では０だが宇宙だけだと誤差が生じる。これが不確定性原理であり、
  異次元が誤差になり、宇宙と加群すると０になる。
  　$ \delta(p,x) $が隠れた変数となり、異次元で誤差をこの不確定性原理
  と表している。
  全ては素粒子方程式から生成される式達である。
   $$ (u + d) + c = e^{-f} - e^{f}, b + (w + s) = e^{-f} + e^{f} $$
   $$ \int \int{1 \over x^s}\mathrm{dvol} - \int \log x \mathrm{dvol}
   = {d \over df}\int\int{1 \over (y \log y)^{1 \over 2}}dy_m
   - \int C dx_m $$
   $$ = e^{-f} - e^{f} $$
   $$ {d \over df}F = {d \over df}\int \int{1 \over (x \log x)^2}dx_m 
   + {d \over df}\int\int{1 \over (y \log y)^{1 \over 2}}dy_m $$
   $$ = e^{f} + e^{-f} $$

     素粒子方程式は、２種類の素粒子と１種類の素粒子の組み合わせで、
   Jones多項式をペアーで構成されている。Rico level理論をベースに
   して、強い力と弱い力、電磁気力が、誤差をD-brane間を重力が
   gravitonとして伝わっているのが、洩れて、反重力をこの
   Weak electric theoryに加えると、全部の力が合わさって、
   統一場理論が完成される。
   これが、ヒッグス場とオイラーの定数を多様体を次元として、
   加群すると、ゼータ関数になる。これが、微分幾何の量子化とガンマ関数
   でもある。

   $$ ||ds^2|| = e^{-2 \pi T|\psi|}[\eta + \bar{h}(x)]dx^{\mu}dx^{\nu}
   + T^2 d^2 \psi $$
   
      素数分布論が、ゼータ関数を開集合として同型でもあり、ノルム空間、経路積分
   として、位相差が原子間力にもなっている。このエネルギー分布が、
   素数と素粒子方程式、宇宙と異次元のフェルミオンとして、一致している。
   このエネルギー分布が、複素多様体である。
   $$\square = 2 (\cos (i x \log x) + i \sin(i x \log x) $$ 


      宇宙と異次元が、双対性として、超対称性の素粒子を形成している。　


 $$ {\partial \over \partial f}F(x) = \int \int \text{cohom}D_k(x)[I_m]  $$
 Cohomology element is constructed with isotopy of D-brane,
 and this isotopy of symmetry structure is sheaf of manifold,
 more over this brane of element is double integrate of projection.
 $$ {\partial \over \partial f}F = {{}^t {\scalebox{1}[2]{{\variint}}}}
 \text{cohom}D_k(x)^{\ll p} $$
 And global partial defferential equation is integrate of cut in cohomolgy,
 and this step of defferential D-brane of sheap value is varint of equals,
 inverse of horizen of cute of section value is reverse of integrate of 
 operators. Three dimension of varint of manifolds in operator of sheaps,
 and this step of times of value is equal with D-brane of equations.
 Global differential equation and integrate equation is operator of 
 global scales of horizen cut and add of cut equation are routs.

 $$ \oint r dx_m = \mathcal{O}(x,y) $$

 $$ {{}^t {\scalebox{1}[2]{{\variiint}}}}_{{D(\chi,x)}} \text{Hom}[D^2 \psi]
 ^{\ll p} \cong  \text{vol}\left({V \over S}\right) $$

 $$ {\partial \over \partial f}F(x) = {{}^t \scalebox{1}[2]{\variint}} 
 \text{cohom}D_k(x)^{\ll p} $$

 $$ {d \over df}F = (F)^{f^{'}}, \int F dx_m = (F)^{f}  $$
 $$ {d \over df}\int \int{1 \over (x \log x)^2}dx_m =
 \left(\int \int{1 \over (x \log x)^2}dx_m \right)^{f^{'}}  $$
 $$ f^{'} = (2x (\log x + \log (x + 1))) $$
  $$  \int \int F dx_m = \left({1 \over (x \log x)^2}\right)^{
	  (\int 2x (\log x + \log (x + 1))dx)} $$ 
   $$= e^{-f} $$
  $$ {d \over df}F = \left({1 \over (x \log x)^2}\right)^{
	  (2x (\log x + \log (x + 1))) } $$  
  $$= e^{f} $$
  $$ \log (x \log x) \ge 2(y \log y)^{1 \over 2} $$

  $$ \log (x \log x) \ge 2(\sqrt{y \log y}) $$

  Global differential manifold is calucrate with x of x times in non
  entropy of defferential values, and out of range in differential formula
  be able to oneselves of global differential compute system.
  Global integrate manifold is also calucrate with step of resulted with
  loop of integrate systems. Global differential manifold is newton and 
  laiptitz formula of out of range in deprivate of compute systems.
  This system of calcurate formula is resulted for reverse of global equation
  each differential and intergrate in non entropy compute resulted values.

  $$ \pi(\chi,x) = [i\pi(\chi,x),f(x)] $$
  $$ \int{1 \over (x \log x)}dx = i \int{1 \over (x \log x)}dx^{N}  
  + {}^{N} i \int{1 \over (x \log x)}dx $$
  $$  \int{1 \over (x \log x)}dx = i \int{x \log x}dx 
  - \int{1 \over (x \log x)}dx $$
  $$ \int \int{1 \over (x \log x)^2}dx_m = i {1 \over 2}x^2 $$
  $$ \int \int{1 \over (x \log x)^2}dx_m = i \int \int_M dx_m $$
  $$ \le {1 \over 2}i + x^2 $$
  $$ E = - {1 \over 2}mv^2 + mc^2 $$
  $$ \lim_{x \to \infty} \int \int{1 \over (x \log x)^2}dx_m \ge {1 \over 2}i $$
  $$ {d \over df}\int \int{1 \over (x \log x)^2}dx_m = {1 \over 2}i $$
  $$ {d \over df}\int \int{1 \over (x \log x)^2}dx_m = 
  \left(\int \int {1 \over (x \log x)^2}dx_m \right)^{(f)^{'}}$$
  $$ (f)^{(f)^{'}} $$
  $$ = (f^f)^{'} $$
  $$ = e^{x \log x} $$

  $$ \Gamma = \int {e^{-x}x^{1-t}}dx, 
  {d \over df}F = e^f, \int F dx_m = e^{-f} $$
  $$ \int x^{1-t}dx = {d \over df}F, \int e^{-x}dx = \int F dx_m $$
  $$ e^{i \theta} = \cos \theta + i \sin \theta $$
  $$ x \bot y \to x \cdot y = \vec{0}, 
  {d \over df}\int \int{1 \over (x \log x)^2}dx_m = {1 \over 2}i $$

  $$ {d \over df}\int \int{1 \over (x \log x)^2}dx_m = 
  \left(\int \int {1 \over (x \log x)^2}dx_m \right)^{(f)^{'}}$$
  $$ {1 \over 2} + {1 \over 2}i = \vec{0}, {1 \over 2}\cdot{1 \over 2}i $$
  $$ A + B  =  \vec{0}, A \cdot B = {1 \over 4}i $$
  $$  \tan 90 \neq 0, ||ds^2|| = A + B $$
  Represented real pole is not only one of pole but also real pole of 
  $ {1 \over 2} $, 
  $\sin 0 = 0$ $$ e^{i \theta} = \cos \theta + i \sin \theta $$
  $\theta $
  this result equation is non-certain system concerned with particle, electric 
  particle and atoms in open set group with possiblity of surface values,
  and have abled to proof in this group of system. Universe and the other
  dimension of concerned of relationship values.
  $$ {d \over df}F = {d \over df} \int \int{1 \over (x \log x)^2}dx_m 
  + {d \over df} \int \int {1 \over (y \log y)^{1 \over 2}}dy_m
  = \Gamma(x,y) $$
  $$ {d \over d \gamma}\Gamma = (e^{-x}x^{1-t})^{\gamma^{'}} \to
  {d \over d \gamma}\Gamma = \Gamma^{(\gamma)^{'}}$$



    $$ \Gamma(s) = \int e^{-x} x^{s - 1}dx, \Gamma^{'}(s)
    = \int e^{-x} x^{s - 1}\log x dx $$
    $$ x^{s - 1} = e^{(s - 1)\log x} $$
    $$ {\partial \over \partial s}x^{s - 1} =
       {\partial \over \partial s}e^{-x}\cdot (e^{{s - 1}\log x}) 
       = e^{-x}\cdot \log x \cdot e^{(s - 1) \log x} $$
    $$ = e^{-x}\cdot \log x \cdot x^{s - 1} $$
    $$ {d \over d \gamma}\Gamma = \Gamma^{(\gamma)^{'}} $$
    $$ {d \over d \gamma}\Gamma(s) = \left( \int^{\infty}_{0}e^{-x}
    x^{s - 1}dx \right)^{\left(\int^{\infty}_{0}e^{-x}x^{s - 1}\log x dx
    \right)^{'}} $$
    $$ = \Gamma^{(\Gamma \int \log x dx)^{'}} $$
    $$ = e^{-x \log x} $$

    $$ \Gamma(s) = \int^{\infty}_{0} e^{-x}x^{s - 1}dx $$
    $$  \Gamma^{'}(s) = \int^{\infty}_{0} e^{-x}x^{s - 1} \log x dx $$
    $$ {d \over df}F = \int x^{s - 1}dx $$
    $$ \int F dx_m = \int e^{-x}dx $$
    $$ {d \over df}F = F^{(f)^{'}}, \int F dx_m = F^{(f)} $$
    Global differential manifold and global integrate manifold are
    reached with begin estimate of gamma and beta function leaded solved,
    Global differential manifold is emerged with gamma function of
    themselves of first value of global differential manifold of gamma
    function, this equation of quota in newton and dalarnverles of equation
    is also of same results equations. Zeta function also resulted with
    quantum group reach with quanto algebra equation.
        Included of diverse of gravity of influent in quantum physics,
    also quantum physics doesn't exist Gauss surface of theorem,
   
    $$ H \Psi = \bigoplus (i \hbar ^{\nabla})^{\oplus L} $$
    $$ = \bigoplus {{ H \Psi} \over \nabla L} $$
    $$ = e^{x \log x} = x^{(x)^{'}} $$
    however,
    $$ {d \over df}F = m(x) $$
    and gravity equation in being abled to existed with quantum physics,
    and this physics is also represented with quantum level of differential
    geometry.
    $$ {d \over dt}\psi(t) = \hbar $$
    $$ = {1 \over 2}i e^{i \Hat{H}} $$
    $$ (i \hbar)^{'} = (- e^{i \Hat{H}})^{'} $$
    $$ = - i e^{i \Hat{H}} $$
    $$ \psi(x) = e^{-i \Hat{H}t}, \bigoplus (i \hbar ^{\nabla})^{\oplus L} 
    = {{{1 \over 2} e^{i \Hat{H}}}^{(-i e^{i \Hat{H}})}} $$
    $$ = ({1 \over 2}f)^{-i f}$$
    $$ = ({1 \over 2})^{-i f} \cdot e^{-x \log x} \cdot (f)^{i} $$
    $$ = \int e^{- x} x^{t - 1}dx, {d \over d \gamma}\Gamma = e^{-x \log x}  $$
     
    $$ f = x, i = t, {1 \over 2} = a, \bigoplus a^{-t x}x^t [I_m]
     \cong \int e^{-x} x^{t - 1}dx $$
     Quantum level of differential geomerty is also resulted with
    Gamma function of global differential manifold of values, Heisenberg
    equation is alos resulted with global integrate manifold and this
    quantum level of differential geometry is manifold of deprive of formula.
    $$ |\psi(t)\rangle_s = e^{-i\Hat{H} t}|\Psi\rangle_{H},
    \Hat{A}_s = \Hat{A}_H(0) $$
    $$ |\Psi(t)\rangle_s \to {d \over dt} $$
    $$ i {d \over dt}|\psi (t) \rangle_s = \Hat{H}|\psi(t)\rangle_s $$
    $$ \langle\Hat{A}(t)\rangle = \langle\Psi(t)|\Hat{A}(0)|\Psi(t)\rangle  $$
    $$ {d \over dt}\Hat{A} = {1 \over i}[\Hat{A},H] $$
    $$ \Hat{A}(t) = e^{i\Hat{H}t}\Hat{A}(0)e^{-i\Hat{H}t} $$
    $$ \lim_{\theta \to 0} \begin{pmatrix} \sin \theta \\
    \cos \theta \end{pmatrix} \begin{pmatrix} \theta & 1 \\
    1 & \theta \end{pmatrix} \begin{pmatrix} \cos \theta \\
    \sin \theta \end{pmatrix} = \begin{pmatrix} 1 & 0 \\
    0 & - 1 \end{pmatrix} $$
    $$ f^{-1}(x) x f(x) = I^{'}_m, I^{'}_m = [1,0] \times [0,1] $$

    $$ {x + y} \ge \sqrt{x y} $$
    $$ {x^{{1 \over 2} + iy} \over e^{x \log x}} = 1 $$
    $$ \mathcal{O}(x) = \nabla_i \nabla_j \int e^{{2 \over m}\sin \theta 
    \cos \theta } \times {{N \mathrm{mod}(e^{x \log x})} \over {O(x)(x +
    \Delta |f|^2)^{1 \over 2}}} $$
    $$ x \Gamma(x) = 2 \int|\sin 2 \theta|^2 d \theta, \mathcal{O}(x)
    = m(x)[D^2 \psi] $$
    $$ i^2 = (0,1) \cdot (0,1), |a||b|\cos \theta = - 1 $$
    $$ E = \mathrm{div}(E, E_1) $$
    $$ \left({\{f,g\}} \over [f,g] \right) = i^2, E = mc^2, I^{'} = i^2 $$
     Gamma function of global differential equation is first of differential
    function deprivate with ordinary differential equations, more also
    universe and other dimensiion of Higgs field emerge with zeta function,
    more spectrum focus is three dimension of manifold of singularity or pair
    theorem, and universe and other dimension of each value equals.
    This reverse of circle function of manifolds is also Higgs fields of 
    dense of entropy and result equation. And differential structure of quantum
    level in gravity is poled with zeta function and quantum group in high 
    result values. Zeta funtion and quantum group is Global differential
    manifold of result in values.
    $$ {d \over d \gamma} \Gamma = m(x) $$
    $$ = e^{- x \log x} $$
    $$ \sin i x = {{e^{-x} + e^{x}} \over 2i} $$
    $$ {d \over df}F = m(x) = e^{f} + e^{-f} $$
    $$ = 2 i \sin (i {x \log x}) $$

    Hyper circle function is constructed with sheap of manifold in 
    Beta function of reverse system emerged from Gamma fucntion 
    and reverse of circle function of entropy value,
    and this reverse of hyper circle function of beta system
    is universe and the other dimension of zeta function,
    more over spectrum focus is this beta function of reverse in
    circle function is quantum level of differential structure oneselves,
    and this quantum level is also gamma function themselves.
    This gamma fucntion is more also D-brane and anti-D-brane built 
    with supplicant values. Higgs fileds is more also pair of universe and
    other dimension each other value elements. 


    My son in acasicrecord demonstrate with being have with quantum level
    of differential structure equation.
    I have with same equation already from being on 
    my father and my doctors of professors give me hints with 
    this qlds equation explanade with global differential manifold.

    $$ H\Psi = \bigoplus {(i \hbar^{\nabla})}^{\oplus L}  (1)$$
    $$ = i\hbar\psi $$
    $$ {d \over df}F = m(x) (2) $$
    $$ = {d \over df}\int \int{1 \over (x \log x)^2}dx_m
    + {d \over df}\int \int{1 \over (y \log y)^{1 \over 2}}dy_m  (3)$$
    $$  \bigoplus {(i \hbar^{\nabla})}^{\oplus L} = \int e^{-x}x^{1 - t}dx_m 
    (4)$$
    $$ = {d \over d \gamma}\Gamma = \int \Gamma(\gamma)^{'}dx_m (5)$$
    $$ = e^f + e^{-f} (6)$$
    $$ {d \over d \gamma}\Gamma^{-1} = e^{f} - e^{-f} (7)$$
    (1) is my son in acasicrecord with his idea, (2),(3) are myself equation.
    (4),(5) are my son and me with tagge of ideas equation.
    (6),(7) are Takashima Aya and me, and my son in acasicrecode with tolio
    equations.
    These references from Grisha professor and Takeuchi Kaoru.
    I read with this references being hint from these professors.
 
      [reference Lisa Randall, Toshihide Masuoka, Makoto Kobayashi]

 
       超重力理論が、超弦理論と一般相対性理論の多様体積分へと、分岐している。
       この超弦理論と一般相対性理論の多様体積分が、超重力理論へと統合される。
        $\theta = 45$度は、平行四辺形の一般座標空間としての変形をすると、
	$\theta = 60$度へと、線形変換することで、$\sin \theta = {\sqrt{3} \over 2}$
	となる。循環となると、永遠の寿命もだめになり、木星の大善は、難しいのは、一目瞭然
	であり、仏教の前後関係でもある。


	These equation and theorem are product with number,
	this element histroy is
	$$ \lim_{n=0}^{\infty}{(1 + {1 \over n})}^{n} = e $$
	$$ C = {{\lim_{x\to 2}}e^{x\log x} \over \log 2}, \int C dx_m = n, 
	 \int e^{-x^2 - y^2}dxdy = \pi $$
	$$ S = \pi r^2, V = \pi ||\int \sin 2x dx||^{2}dx, 
	 V - \int \sin x dx = \beta(p,q), \beta(p,q) = \bigoplus
	{(i \hbar^{\nabla})}^{\oplus L} = {d \over {d f}}F(x,y) $$
	$$ = \int \Gamma^{'}(\gamma)dx_m 
          = (x,y)\cdot (x,y) = e^{ix} $$
	$$ e^{ix} = \cos x + i \sin x 
	 = {d \over {d f}}e^{ix} = \sin(i\log x) + i\cos(i\log x) $$
	$$ = n $$
	This develop and limitation element escourt into Euler product 
	emerged from general of Euler equation.
	And, this equation built with Higgs field and zeta function
	from gamma function of partial integral deprivation manifold.

	$$ \pi \int ||[\Phi / \pi r^2]||dr, {}^{t}\scalebox{1.0}[2]{
		\variint} \Delta(\pi(\chi,x))[I_m], 
		 \sum(\sigma(H \times K)) $$
		$$ [\alearletter / \nabla]^{\mu\nu}, 
         \secproduct \otimes \Delta $$
	 Secure product construct with anti-gravity equation integrate with
	 average equation.
	$$ \nabla_{i}\nabla_{j}(\square \times \alearletter)d\tau, 
	 \sqrt{x_m\cdot y_m} $$
	$$ \squareiint \mathrm{cohom}D_{\chi}[I_m],
	\bigotimes [S_{D\chi} \otimes h\nu] $$
	$$ \ll i\hbar \psi \ast H\Psi \gg, 
	 \int \Delta(\zeta)d\zeta $$

	 Step function emerge with Thurston Perelman manifold.
	$$ \whitewithbodercircleint (I_m)^{\nabla L} $$
	$$ -2\int{{\nabla_{i}\nabla_{j} (R + \nabla\nabla f)} \over
	{\Delta(R + \Delta)}}dm $$
	$$ \Delta(F(\Delta) \times \Delta(G(\Delta))) = - (F(\Delta) \cup
	F(\Delta)) + (F(\Delta) \cap F(\Delta)) $$
	$$ \sum \square(\nabla)[I_m], \secproduct,[\nabla / \square],
	\pi ||\int \nabla_{i}\nabla_{j} f \nabla d \eta||^2 = S^{m}
	\times S^{m-1} $$
	$$ (\pi r^2) dr_m, {{\partial^{df} (e^{f} + e^{-f})} \over {
		d(e^{f} - e^{-f})}} $$
	$$ \zeta = (at -t^n + a) = e^{f} $$
	$$ \text{Euler product is} C = 2e, C ~ {e \over 5} $$

	Rotate of pole on imaginary and real estimation construct with
	non commutative manfiold, and this equation Quantum physics
	create from zeta function.
  $$ \nabla \Psi \otimes \phi D = \Phi [M^{\nabla}]^{\smallaverage} $$
  $$ {}^{t}\variiint Dcohol[I_m] $$
  $$ {}^{t}\variint \chi(\chi,d)d\chi $$
  $$ {}^{t}\variint [i\pi(\chi,x)]d\chi $$
  Dalanversian equation rotate with cohomolog element from being Heisenberg
  equation emerged into fundarment group belong with quato of empty set.
  $$ \square^{\ll p}\flowerletter cohom D\chi[I_m] $$
  $$ = \Psi\bigotimes i\hbar\psi $$
  $$ H\Psi \bigotimes i\hbar\psi = \square^{\ll p} $$
  $$ {}^{t}\scalebox{1.0}[2]{\variint}_{D} 
  \pi(\chi,x)_m = \ll |\emptyset | \alearletter| \gg $$

 Therefore, quantum calcuration systemiate with Euler product create into
 Jones manifold.
  $$ \scalebox{1.0}[2]{$\times$} ||ds^2|| = \ll \square | \alearletter \gg $$
  $$ {{\squareiint N} \over {\emptyset M}}d\tau  = {{||ds^2||} \over {\emptyset
   {M}}} = {e^{x\log x} \over {\log 2}}, \int{{ [\pi r^2]^{\circ i}} \over 
   {\emptyset M}} = e^{f} + e^{-f} $$
  

\end{document}
